\part{Sünde und Vergebung}
\chapter{Verunreinigung des Landes}
tl;dr:
Manche Sünden sind so groß, dass sie nicht bloß das Heiligtum verunreinigen, sondern dass sie `das Land' selbst verunreinigen.\footnote{
Das sind: Götzendienst, Sexuelle Vergehen, und Blutvergießen.}
Für diese Sünden gibt es keine Reparationsopfer und auch ihre Verunreinigung kann nicht einmal durch Säuberungsopfer oder Yom Kippur entfernt werden, weil
diese Opfer nur das Heiligtum dekontaminieren.
Die einzigen Möglichkeiten, mit diesen Opfern umzugehen ist, die Todesstrafe zu vollziehen.
Wenn das Land zu sehr verunreinigt worden ist, folgt das Exil, damit das Land Ruhe hat und durch diese verstrichene Zeit wieder rein werden kann (in Analogie zur Reinigung
der meisten rituellen Unreinheiten die unter anderem durch Zeit geschieht).

\chapter{\label{chap:bundesfluch}Der Lohn der Sünde und der Fluch des Bundes}
tl;dr:
`Der Lohn der Sünde ist der Tod' und das bedeutet, dass Sünde unweigerlich zum Tod führt.
Tod ist nicht immer eine Strafe Gottes.
Manchmal tötet Gott Menschen als Strafe für ihre Sünden, aber nicht jeder Tod ist Gottes Strafe; Tod ist die natürliche Folge der Sünde.
Tod ist `die Macht des Teufels', und `der letzte Feind, der vernichtet wird'.
Die Sünde verwendet das perfekte Gesetz, das zum Leben führen sollte, um stattdessen Tod hervorzurufen.


% TODO Frage:
% Ist Die "Strafe des Gesetzes" ausschlißlich eine natürliche Konsequenz, genau so wie Tod eine natürliche Konsequenz der Sünde ist?
% - insb. interessant für Gal 3:13
% Die Idee wäre dann:
% "Das Gesetzt bringt leben dem der es befolgt. Durch die Sünde sind Menschen nicht in der Lage es zu halten. Als Konsequenz für das Nicht-Halten kommt Fluch. Jesus besiegt die Sünde, indem er die Konsequenz des Gesetztes trägt, uns reinwäscht und seine immortalität anzieht, um uns von der Sünde zu heilen."
% 
\chapter{\label{chap-fiat}Vergebung durch Fiat}
tl;dr:
Gott vergibt einfach so - meist nachdem die Person Buße getan hat.
Gott muss Sünde nicht bestrafen.
Jesus vergibt in seinem Leben einfach so.
Nie hat sich jemand vor Jesu Tod gefragt wie es möglich ist, dass Gott Sünden vergeben kann ohne sie zu bestrafen.
Gott entscheidet sich durch Fiat, Sünden zu vergeben.

\chapter{\label{chap:prophetische-hoffnung}Die prophetische Hoffnung}
tl;dr:
Die Propheten leben kurz vor oder während des Exils, also in der Zeit in der Israel unter dem Fluch des Gesetzes steht.
Ihre Hoffnung ist nicht, dass es ein endgültiges, finales Säuberungsopfer geben wird, mit dem Gott alle Sünden abwaschen kann.
Stattdessen ist die Prophetische Hoffnung konsistent so beschrieben:
(a) Gott wird Israel aus Gnade wieder annehmen und ihnen vergeben
(b) Gott wird die Verunreinigung der Sünden durch eine Waschung im Heiligen Geist von den Menschen entfernen
(c) Gott wird die Menschen grundsätzlich verändern, um sie zu befähigen, Gott zu dienen.
Weiterhin weiten die Propheten die Kategorie der Sünden, die das Land verunreinigen auf, indem sie Ökonomische Unterdrückung als eine Form von Blutvergießen beschreiben.

% Israel noch (teilweise / spirituell) im Exil (2. Tempel Texte zitieren)
% 4Q504 5:15-16; Isa 44:1-3 :: Relation zwischen dem Ende des Exils und der Ausgießung des heiligen Geistes

\chapter{Der Messias und die Kräfte des Todes}
tl;dr:\footnote{
Dieses Kapitel ist eine kurze Zusammenfassung von \Cite{jesus-and-the-forces-of-death}}
Jesus tritt in den Evangelien auf und kommt immer wieder in Kontakt mit den verschiedenen Formen ritueller Unreinheit.
Anstatt die Menschen nur rituell zu reinigen, heilt er stattdessen die Quelle der Unreinheit - den Aussatz, den Blutfluß, den Tod an sich.
Alle Quellen der rituellen Unreinheit wurden als Kräfte des Todes aufgefasst, und überall wo der Messias mit den Kräften des Todes in Kontakt kommt besiegt er
sie durch sein überfließendes Leben und seine Heiligkeit.
Das große Finale seines Wirkens erwirkt Jesus durch seinen Tod und seine Auferstehung.
Hier kommt er mit dem Tod an sich in Kontakt, und wieder ist sein göttliches Leben dem Tod überlegen - Jesus vernichtet den Tod durch sein Leben.

\chapter{Der Tod des Messias und der neue Bund}
tl;dr:
Jesus selbst setzt das Abendmahl als eines der zentralen Rituale für seine Nachfolger ein.
Dieses Mahl ist explizit Jesus als neues Passahlamm, und die Erinnerung an einen neuen Bundesschluss.
Jesus vergießt sein Blut um diesen Bund einzuleiten und seine Nachfolger dadurch zu reinigen und zu heiligen.
Jesu Wirken erfüllt genau die prophetische Hoffnung, und beendet das Exil endgültig.

\chapter{\label{chap-hebrews}Der Messias und sein Opfer im Hebräerbrief}
Nachdem die Evangelien Jesu Tod als Einläuten des neuen Bundes, die Erfüllung der prophetischen Hoffnung, und seinen Sieg über die Kräfte des Todes dargestellt haben wenden wir uns jetzt dem Hebräerbrief zu.
Hier haben wir den längsten Abschnitt des neuen Testaments, der sich in detaillierter Art und Weise mit der Bedeutung von Jesu Opfer in Bezug auf die Opfer im levitischen System auseinandersetzt.
Es ist als allererstes sehr klar, dass Jesus sich selbst als Opfer dem Vater darbringt, und dass dieses Opfer dafür sorgt, dass Menschen Vergebung
für ihre Sünden erhalten. Dieser Punkt steht nicht zur Debatte.
Wenn wir den Hebräerbrief und seine Interaktion mit Jesus als Opfer verstehen, dann wird uns das ein klares Bild davon geben, wie PSA nicht nötig ist
um mit vollem Herzen sagen zu können `Jesus ist das Opferlamm, das zur Vergebung unserer Sünden geschlachtet wurde.'

Im gesamten Hebräerbrief legt der Author dar, wie Jesus das Alte Testament erfüllt, bloß in größerer, besserer, vollständigerer Form.
Da wir uns speziell mit dem levitischen Opfersystem und Jesus als Opfer auseinandersetzen beginnt die interessante Diskussion für unsere Frage in Kapitel 7.
Hier begründet der Author (a) warum Jesus ein Priester sein kann, obwohl er nicht von Levi abstammt (b) warum er sogar besser ist als die irdischen Priester.
Beides wird begründet, da Jesus ein Hoher Priester in der Ordnung Melchisedeks ist - dem prototypischen Priester Gottes außerhalb des jüdischen Volkes.
Kapitel 8 führt das Argument weiter: Wenn Jesus ein Hoher Priester ist, der besser ist als die irdischen Priester, dann setzt er einen Dienst und einen Bund ein, der besser ist.
Aus diesem Gedanken des neuen Bundes heraus wird die gesamte Diskussion in Kapiteln 9 bis 10 fließen.
Der neue Bund ist besser, und seine Opfer und Riten sind besser.
Dieser neue Bund wird direkt mit der prophetischen Hoffnung in \bibverse{Jeremiah}{31}{} verbunden, das der Autor wortreich zitiert.
Dieser neue Bund wird `das Gesetz in ihren Sinn geben'\bibfoot{Hebrews}{8}{10}, und verbunden mit dem Schluss des neuen Bundes wird Gott `ihrer Sünden nicht mehr gedenken.'\bibfoot{Hebrews}{8}{12}.

Kapitel 9 bis 10 füllen jetzt die Details dieses Vergleiches aus. Wenn Jesus Hoher Priester eines neuen Bundes ist, was sind die Opfer und Riten dieses Bundes? Wie funktioniert er, und was bewirkt er?
Diese Kapitel werden wir jetzt Punkt für Punkt durchgehen und insbesondere sehen, wie Jesus als Opfer beschrieben wird.
Verse 11 bis 12 beginnen wie folgt:
\begin{quote}
Christus aber ist gekommen als Hoher Priester der zukünftigen Güter und ist durch das größere und vollkommenere Zelt –
das nicht mit Händen gemacht, das heißt, nicht von dieser Schöpfung ist – und nicht mit Blut von Böcken und Kälbern,
sondern mit seinem eigenen Blut ein für alle Mal in das Heiligtum hineingegangen und hat \textlangle uns\textrangle\ eine ewige Erlösung erworben
\attrib{\bibverse{Hebrews}{9}{11-12}, \Cite{elb}}
\end{quote}
Jesus ist besser als der alte Bund, und dieses besser beginnt damit, dass sein Blut besser ist als das von Tieren.
Wir erinnern uns an \bibverse{Leviticus}{17}{11} und die Diskussionen darüber, dass das Blut der Opfertiere oft der primäre Grund ist, sie zu Opfern.
Jesus Blut ist sein Leben, mehr als sein Tod.
Und mit seinem Leben (rituell angewandt, ausgeschüttet am Kreuz) tritt Jesus in das himmlische Heiligtum hinein.
Jesus \textit{ist} einfach von Natur aus der einzige Mensch, der aus sich selbst hinaus in das Heiligtum herauf steigen kann.\footnote{
Athanasius in \Cite{de-incarnatione-verbi} ist besonders fokussiert darauf, dass Jesus durch seine Inkarnation die göttliche Natur mit der menschlichen teilt,
was sehr gut mit der Feststellung harmonisiert, dass Jesus durch seine Natur das vollständige Leben ist, dass den Tod überwindet.
Wir werden in unserer Diskussion über die Kirchenväter dazu zurückkehren: Soteriologie ist Christologie.}

\begin{quote}
Denn wenn das Blut von Böcken und Stieren und die Asche einer jungen Kuh, auf die Unreinen gesprengt, zur Reinheit des Fleisches heiligt,
wie viel mehr wird das Blut des Christus, der sich selbst durch den ewigen Geist \textlangle als Opfer\textrangle\ ohne Fehler Gott dargebracht hat,
euer Gewissen zu reinigen von toten Werken, damit ihr dem lebendigen Gott dient!
\attrib{\bibverse{Hebrews}{9}{13-14}, \Cite{elb}}
\end{quote}
Hier ist das erste Ritual, das Jesus durch sein besseres Blut möglich macht: Jesus wäscht uns rein und heiligt uns.
Wir haben hier zwei Funktionen mit zwei verschiedenen Typen im Alten Bund:
(a) die körperliche Reinigung durch die Asche einer jugen Kuh bezieht sich auf \bibverse{Numbers}{19}{} und das Reinigungsritual für Leichen-Unreinheit.
So wie dieses physische Reinigungsmittel die rituelle Unreinheit von einem Menschen entfernt, der eine Leiche berührt hat, wäscht Jesu Blut unsere ``toten Werke''
ab.
(b) Das Blut von Böcken und Stieren auf die Unreinen gesprengt bezieht sich wohl auf die Anwendung von Blut beim Bundesschluss oder bei der Priesterweihe.\footnote{
	Die Blutanwendungen bei der Reinigung von Aussatz ist unwahrscheinlich, da das Blut dort auf den ehemals Aussätzigen `gegeben' wird, nicht `gesprengt'.
	Vergleiche \bibverse{Leviticus}{13}{}.
}
In \ref{metaphysische-transformation} haben wir bereits untersucht, was diese Besprengung bewirkt: eine Transformation hin zu einer größeren Heiligkeit.
Genau so ist es Jesu Blut, dass die Gläubigen reinigt und transformiert: es wäscht unsere toten Werke ab (wie die Asche einer jungen Kuh die Leichen-Unreinheit abwäscht),
und es macht uns heilig, wie das Blut gesprengt auf die werdenden Priester sie für ihre Aufgabe in Gottes direkter Gegenwart heiligt.
Das dekontaminierende Opfer hier nicht gemeint sein können sehen wir sofort daran, dass ihr Blut niemals auf Menschen angewandt wurde.
\footnote{Man könnte einwenden, dass προσήνεγκεν hier Jesus als Opfer darstellen muss und die Bedeutung als Reinigungsmittel ausschließt.
	Doch obwohl `opfern' eine natürlich Bedeutung ist, taucht προσφέρω in LXX genau so für nicht-sakrifiziale (in der Tat: nicht-rituelle) Handlungen auf.
	Die Einfügung \textlangle als Opfer\textrangle\ der \Cite{elb} ist hier nicht zwingend nötig, und dieses Verb kann beide Handlungen einschließen.
}
Jesu Blut ist das perfekte unausschöpfliche Leben par excellence - das beste Reinigungsmittel der Welt.
Und so reinigt es unser Gewissen von toten Werken, wo die Reinigung im levitischen System nur den Körper von der Unreinheit des Todes reinigen konnte.
Wir dienen einem lebendigen Gott, reingewaschen von unseren toten Werken und wir tun es durch die Waschung in Jesu göttlichem Leben -
hier sehen wir sofort die Antitypen der beiden Zentralen Punkte der prophetischen Hoffnung:
Waschung durch Gott, und Befähigung ihm zu dienen\footnote{
Denke an \bibverse{Jeremiah}{31}{34}, das der Autor gerade erst zitiert hat.
}.
Mit diesem neuen Bund schließt auch Verse 15 direkt an:
\begin{quote}
	Und darum ist er Mittler eines neuen Bundes, damit, da der Tod geschehen ist zur Erlösung von den Übertretungen unter dem ersten Bund, die Berufenen die Verheißung des ewigen Erbes empfangen.
	\attrib{\bibverse{Hebrews}{9}{15}, \Cite{elb}}
\end{quote}
Jesu Tod bringt `Erlösung der Übertretungen unter dem ersten Bund', und nicht nur das - er ermöglicht `den Berufenen die Verheißung des ewigen Erbes' zu empfangen.
Beide diese Konzepte sind direkt mit dem neuen Bundesschluss verbunden:
Der neue Bund ist es, der Vergebung der Sünden mit sich bringt nachdem unsere Natur und Sünde gewaschen wurde, genau wie Jeremiah vorher gesagt hatte.
Vergebung und Erlösung ist aus Gottes Gnade allein, durch den neuen Bund versiegelt im Opfer des Messias, für alle die an diesem Bund teilhaben.
Aber warum war es nötig, dass Jesus gestorben ist?
Hätte der neue Bund nicht ohne ein Opfer besiegelt werden können?
Zu dieser Frage wendet sich der Autor in den nächsten Versen:
\begin{quote}
	Denn wo ein Testament ist, da muss notwendig der Tod dessen eintreten, der das Testament gemacht hat. Denn ein Testament ist gültig, wenn der Tod eingetreten ist, weil es niemals Kraft hat, solange der lebt, der das Testament gemacht hat.
	\attrib{\bibverse{Hebrews}{9}{16-17}, \Cite{elb}}
\end{quote}
Interessant ist hier vor allem, dass der Autor mit der konsistenten Metapher- und Typ-Antityp-Beziehung zum levitischen System bricht und plötzlich eine legale Metapher verwendet, um zu zeigen, warum Jesus Tod notwendig war.
Dass er das tut ist aber sehr verständlich wenn wir beachten, dass Opfer im levitischen System so explizit als Nicht-Tode formuliert wurden (wie in Kapitel \ref{opfer-als-nicht-tod} besprochen).
Jesu Opfer war nötig, damit sein Blut zu unserer Waschung zur Verfügung steht.
Sein tatsächlicher Tod war auch nötig - aber um das zu begründen bedient sich der Autor bewusst einer Metapher, die nicht das Opfer direkt mit dem Tod in Verbindung bringt.
Zu der Notwendigkeit des Blutes kehren wir prompt zurück:
\begin{quote}
Daher ist auch der erste ⟨Bund⟩ nicht ohne Blut eingeweiht worden.
Denn als jedes Gebot nach dem Gesetz von Mose dem ganzen Volk mitgeteilt war, nahm er das Blut der Kälber und Böcke mit Wasser und Purpurwolle und Ysop und besprengte sowohl das Buch selbst als auch das ganze Volk
und sprach: »Dies ist das Blut des Bundes, den Gott für euch geboten hat.«
Aber auch das Zelt und alle Gefäße des Dienstes besprengte er ebenso mit dem Blut;
	\attrib{\bibverse{Hebrews}{9}{18-21}, \Cite{elb}}
\end{quote}
Wieder und wieder dreht sich hier alles um Jesu Blut als Blut des Bundes, das selbst im Heiligtum angebracht wurde zur Ratifizierung des Bundes.
Es wäscht, aber vor allem heiligt es, wie wir in Kapitel \ref{metaphysische-transformation} gesehen haben.
Im folgenden möchte der Autor Jesu Werk für die Dekontamination des himmlischen Heiligtums erläutern, und um den Bogen vom Blut des Bundes zum Blut der Dekontamination zu schlagen steht Vers 22:
\begin{quote}
	und fast alle Dinge werden mit Blut gereinigt nach dem Gesetz, und ohne Blutvergießen gibt es keine Vergebung.
	\attrib{\bibverse{Hebrews}{9}{22}, \Cite{elb}}
\end{quote}
Der erste Satz ist leicht verständlich: so wie alle Dinge mit Blut geheiligt werden, werden sie auch dekontaminiert.
Der zweite Teilsatz ist der komplizierteste Vers in diesen zwei Kapiteln des Hebräerbriefes, einfach weil er - naiv gelesen - falsch ist.
Wir haben bereits im Detail beobachtet, dass Vergebung unabhängig von Blut und Opfern geschehen kann (in Kapitel \ref{chap-fiat}).
Außerdem wissen wir, dass das Resultat der Säuberungsopfer die in v22a offensichtlich im Fokus stehen nicht Vergebung, sondern Dekontamination ist (Kapitel \ref{kipper}).
Nicht nur das, sondern selbst wenn Säuberungsopfer Sünden vergeben würden, wird nicht bei allen Säuberungsopfern Blut vergossen.
Hier sind diverse Möglichkeiten, was `ohne Blutvergießen gibt es keine Vergebung' bedeuten kann:
\begin{description}
	\item[Gott kann keine Sünden vergeben, ohne dass Blut fließt] Diese Aussage ist wie schon diskutiert einfach falsch.
	\item[Freilassung durch den neuen Bund] Diese Idee passt zwar zur bisherigen Diskussion in Kapitel 9, aber sie unterbricht den Umschwung der Argumentation von Bundesschluss zu Jesus als Säuberungsopfer und ist daher unwahrscheinlich.
	\item[Freilassung von der Notwendigkeit, Säuberungsopfer zu bringen] Diese Aussage ist mit Sicherheit richtig, erscheint aber recht kompliziert.
	\item[Freilassung von der Dekontamination an sich] Diese Option muss abgelehnt werden weil sie verlangt, dass ἄφεσις mit \kipper\ gleichzusetzen ist - diese Gleichsetzung ist dadurch unwahrscheinlich, dass die LXX niemals ἄφεσις (oder seine verbale Form) für kipper verwendet, sondern nur für נִסְלַח (nislaḥ, vergeben) in diesen Kontexten.
		Wenn der Kontext für die Verwendung von ἄφεσις im levitischen System zu finden ist, dann kann dieses Wort nur schwer Dekontamination bedeuten.
	\item[Vergebung von Sünden für die es Säuberungsopfer gab] Auch diese Aussage erscheint kompliziert, aber zumindest fließt sie gut mit dem Argument.
		(a) Es geht tatsächlich weiter um Säuberungsopfer
		(b) das Blutvergießen wird konsistent mit v22a und v23 als Blut eines dekontaminierenden Opfers gelesen
		(c) diese Aussage ist explizit eine Vorschrift aus dem Gesetz, das in v22 angebracht wird und
		(d) bis auf das Problem der Säuberungsopfer für Arme, die ohne Blut ablaufen, ist die Aussage tatsächlich richtig und wird nicht hier das erste mal so formuliert.\footnote{
		Gott wird Sünden für die es Säuberungsopfer gibt nur vergeben, wenn diese Opfer auch gebracht werden.
	Diese Aussage ist so unproblematisch, dass sie in einem Halbsatz ausgesprochen werden kann ohne vom eigentlichen Argument abzulenken.}
\end{description}
Auch die letzte Option erscheint etwas kompliziert, aber es ist meiner Ansicht nach die beste Interpretation dieses sehr schwieriegen Verses.

Die Argumentation des Briefes schreitet jetzt fort:
\begin{quote}
\textlangle Es ist\textrangle\ nun nötig, dass die Abbilder der in den Himmeln \textlangle befindliche Dinge\textrangle\ hierdurch gereinigt werden, die himmlischen Dinge selbst aber durch bessere Schlachtopfer als diese.
Denn Christus ist nicht hineingegangen in ein mit Händen gemachtes Heiligtum, ein Abbild des wahren \textlangle Heiligtums\textrangle\, sondern in den Himmel selbst,
um jetzt vor dem Angesicht Gottes für uns zu erscheinen, auch nicht, um sich selbst oftmals zu opfern, wie der Hohe Priester alljährlich mit fremdem Blut in das Heiligtum hineingeht
– sonst hätte er oftmals leiden müssen von Grundlegung der Welt an –; jetzt aber ist er einmal in der Vollendung der Zeitalter offenbar geworden, um durch sein Opfer die Sünde aufzuheben.
	\attrib{\bibverse{Hebrews}{9}{23-26}, \Cite{elb}}
\end{quote}
So wie ein irdisches Säuberungsopfer das irdische Heiligtum reinigt, so reinigt Jesu perfektes Blut das himmlische Heiligtum.
Das ist eine Facette der `Versöhnung', die in Christus passiert: alle Unreinheit, die durch Sünde\footnote{und ggf. rituelle Unreinheit, aber ob rituelle Reinheit für den Autor noch eine Rolle spielt ist sehr fragwürdig} auf das himmlische Heiligtum gekommen ist wird gereinigt.
Nur so ist es möglich, dass das himmlische Heiligtum seine Funktion weiterhin erfüllen kann.
Das ist `Versöhnung' durch sein Blut - das Heiligtum ist gereinigt und wir haben Zugang zum Vater, durch Jesu Opfertod am Kreuz.
Das Problem, dass der Hebräerbrief hier löst ist `das himmlische Heiligtum ist verunreinigt' - und die Lösung ist Jesu Opfertod für unsere Sünden.
Substitionelle Bestrafung des perfekten Hohe Priesters des neuen Bundes, der sich selbst als Opfer darbringt ist in dieser Passage nicht im Ansatz indikiert.
Wenn Jesu Tod als stellvertretende Strafe zur Vergebung von Sünden das Ziel dieser Passage wäre, müsste der Autor gleich mit mehreren Tabus des levitischen Systems brechen (Opfer sind keine Strafe; Opfer sind kein stellvertretender Tod; Opfer sind nicht einmal Tod; Opfer haben eine sehr klar definierte andere Funktion), aber das ist nicht was wir sehen.
Der Autor kennt sein Levitikus, und er verwendet das vielschichtige rituelle System der Torah mit großer Vorsicht und Klarheit.
Der Messias und sein Tod als Opfer für unsere Sünden erfüllt alle rituelle Funktionen des levitischen Systems.
Wir sind reingewaschen, wir sind geheiligt als Teilhaber an einen neuen Bund, das Heiligtum ist gereinigt.
Diese bereits etablierten Funktionen werden jetzt durch den Rest der Passage noch weiter ausformuliert, um weiter klar zu machen dass Jesu Opfer größer ist als es die levitischen jemals sein konnten.

Verse 27 und 28 bilden eine Vorausschau auf Jesu zweites Kommen und das endgültige Gericht, und zum Beginn von Kapitel 10 fahren wir mit unserem Thema - Jesus und das levitische System - fort.
\begin{quote}
Denn da das Gesetz einen Schatten der zukünftigen Güter, nicht der Dinge Ebenbild selbst hat, so kann es niemals mit denselben Schlachtopfern,
die sie alljährlich darbringen, die Hinzunahenden für immer vollkommen machen.
Denn würde sonst nicht ihre Darbringung aufgehört haben, weil die den Gottesdienst Übenden, einmal gereinigt, kein Sündenbewusstsein mehr gehabt hätten?
Doch in jenen \textlangle Opfern\textrangle\ ist alljährlich ein Erinnern an die Sünden;
denn unmöglich kann Blut von Stieren und Böcken Sünden wegnehmen.
	\attrib{\bibverse{Hebrews}{10}{1-4}, \Cite{elb}}
\end{quote}
Diese Aussagen sollten uns zum jetzigen Zeitpunkt sehr bekannt vorkommen: Natürlich können die levitischen Opfer keine Sünden wegnehmen - nie war das ihre Funktion.
Dass das eigentliche Problem, nämlich die Sünde an sich, von ihnen nie gelöst wurde sehen wir daran, dass diese Opfer immer wieder dargebracht werden müssen.
Jetzt aber hat Christus ein Opfer gebracht, nicht bloß um das Heiligtum zu reinigen und nicht bloß um Vergebung auszusprechen - sondern um mit diesem Opfer eine wirkliche Transformation seines Bundesvolkes zu bewirken:
\begin{quote}
Darum spricht er, als er in die Welt kommt: »Schlachtopfer und Opfergabe hast du nicht gewollt, einen Leib aber hast du mir bereitet;
an Brandopfern und Sündopfern hast du kein Wohlgefallen gefunden.
Da sprach ich: Siehe, ich komme – in der Buchrolle steht von mir geschrieben –, um deinen Willen, Gott, zu tun.«
Vorher sagt er: »Schlachtopfer und Opfergaben und Brandopfer und Sündopfer hast du nicht gewollt, auch kein Wohlgefallen daran gefunden« –
die doch nach dem Gesetz dargebracht werden –;
dann sprach er: »Siehe, ich komme, um deinen Willen zu tun« – er nimmt das Erste weg, um das Zweite aufzurichten.
In diesem Willen sind wir geheiligt durch das ein für alle Mal geschehene Opfer des Leibes Jesu Christi.
	\attrib{\bibverse{Hebrews}{10}{5-10}, \Cite{elb}}
\end{quote}
Denn das Opfer Jesu bewirkt etwas größeres, als all die Opfer des alten Bundes:
Jesu Opfer `heiligt' uns in Analogie dazu, wie die Opfer des Bundesschlusses das Volk geheiligt haben, aber auf eine viel bessere und größere Art und Weise.\footnote{
\bibverse{Psalms}{40}{}, aus dem hier zitiert wird, enthält eine der beinahe üblichen Aussprüche darüber, dass Gott das Herz wichtiger ist als die Opfer an sich.
Außerdem versteht der Autor des Hebräerbriefes diesen Psalm scheinbar messianisch, und stellt in seinem Midrasch Jesu Wirken und speziell sein freiwilliges Opfer als das da, was Gott wirklich wollte:
Den perfekten Gott-Mensch, dessen Leben und Tod Gott ein wohlgefälliges Opfer ist.}
An den nächsten Versen können wir sehen, dass der Fokus des Autors hier tatsächlich konsistent auf Jesu Opfer als Antitypos des Bundesschlusses -
und insbesondere der metaphysischen Transformation des Bundesvolkes - liegt.
Insbesondere ist also der kurze Ausflug in eine Diskussion über Jesus als Säuberungsopfer bereits beendet:
\begin{quote}
Und jeder Priester steht täglich da, verrichtet den Dienst und bringt oft dieselben Schlachtopfer dar, die niemals Sünden hinwegnehmen können.
Dieser aber hat ein Schlachtopfer für Sünden dargebracht und sich für immer gesetzt zur Rechten Gottes.
Fortan wartet er, bis seine Feinde hingelegt sind als Schemel seiner Füße.
Denn mit einem Opfer hat er die, die geheiligt werden, für immer vollkommen gemacht.
⟨Das⟩ bezeugt uns aber auch der Heilige Geist; denn nachdem er gesagt hat:
»Dies ist der Bund, den ich für sie errichten werde nach jenen Tagen, spricht der Herr, ich werde meine Gesetze in ihre Herzen geben und sie auch in ihren Sinn schreiben«;
und: »Ihrer Sünden und ihrer Gesetzlosigkeiten werde ich nicht mehr gedenken.«
	\attrib{\bibverse{Hebrews}{10}{11-17}, \Cite{elb}}
\end{quote}
Wir sind vielleicht geneigt zu lesen ``Dieser aber hat ein Schlachtopfer für Sünden dargebracht'' und daraus zu lesen, dass `Sein stellvertretendes Sühneopfer unsere Schuld auf sich nimmt', aber eine solche Interpretation ist mit einem klaren Verständnis der Argumentation und der Hintergründe im levitischen System einfach unmöglich.
Um es nocheinmal zu sagen:
Kein Opfer hat jemals `stellvertretend Sünden vergeben/bestraft',
der Autor hat nie argumentiert, dass sich dieser Zustand jetzt ändert,
und wir dürfen den Text nicht in unser falsches Verständnis des levitischen Systems zwingen.
Tatsächlich ist der Text sehr eindeutig darin, welche Opfer hier angesprochen werden und in welcher Art und Weise genau diese Opfer `für Sünden' sind.
Bereits in Vers 10 war das Ziel des Opfers `Heiligung'; in Vers 12 ist es `für Sünden' und in Vers 14 zur `Vervollkommnung'.
Direkt danach, und direkt als Begründung verbunden wird wiederum die Hoffnung des neuen Bundes aus \bibverse{Jeremiah}{31}{} zitiert.
Jesu Opfer in Versen 10 bis 14 ist ganz eindeutig ein Opfer des Bundesschlusses, kein Säuberungsopfer.
Jesu Blut, das am Kreuz für unsere Sünden vergossen wurde läutet den neuen Bund ein, auf den die Propheten gehofft haben.
Es heiligt uns und bringt uns zur Vollkommenheit, indem es uns transformiert.
Ja, unsere Sünden sind uns vergeben - denn so einen Gott haben wir, der Sünder zu seinem Bundesvolk macht und ihnen vergibt.
Aber das Opfer des fleischgewordenen Gottes bewirkt viel mehr als das.
Das perfekte Blut des inkarnierten Wortes, der nicht anders kann als zu leben ist es, durch das wir \textit{leben}, durch den wir \textit{heilig sind} und durch das \textit{die Ursache für die Sünde} zusammen mit der Strafe für diese Sünde ausgelöscht wird: unsere gefallene, korrupte Natur.
Dazu war Christi Tod nötig: um den neuen Bund einzuleiten - und genau das hat der Autor explizit in \bibverse{Hebrews}{9}{16-17} argumentiert.

\begin{quote}
Wo aber Vergebung dieser \textlangle Sünden\textrangle\ ist, gibt es kein Opfer für Sünde mehr.
	\attrib{\bibverse{Hebrews}{10}{18}, \Cite{elb}}
\end{quote}
Hier begegnet uns eine besondere Feinheit, die in der deutschen Übersetzung leider schwer zu transportieren ist:
Der technische Term für die Säuberungsopfer in der LXX ist περὶ ἁμαρτίας (peri hamartias, für Sünden), und diese Konstruktion ist uns in den Zitaten aus Psalm 40 begegnet, wo die Elberfelder Bibel `Sündopfer' übersetzt hat.
Interessanterweise hat der Autor in Vers 12 `ein Schlachtopfer für Sünden' hier mit ὑπὲρ ἁμαρτιῶν προσενέγκας θυσίαν (hyper hamartion prosenegkas thysian) wiedergegeben - einem anderen grichischen Wort das hier im Deutschen legitim `für' übersetzt werden kann.
In Vers 18 kehrt der Autor jetzt zu περὶ ἁμαρτίας zurück - dem griechischen Begriff, der explizit mit Säuberungsopfern verbunden ist.
Wo ein neuer Bund geschaffen ist, um die Sünden zu vergeben und abzuwaschen gibt jetzt keine anderen Säuberungsopfer mehr, um sich um die Sünden zu kümmern.\footnote{
Selbst die Sünden, für die vorher Säuberungsopfer vorgeschrieben waren - alles ist mit einem finalen, perfekten Opfer beseitigt.}
Der neue Bund und das neue Opfer ist so viel besser, dass ``er den ersten für veraltet erklärt''\bibfoot{Hebrews}{8}{13}.

In den folgenden Versen wird das theoretisch Erklärte für die Zuhörer praktisch:
\begin{quote}
Da wir nun, Brüder, durch das Blut Jesu Freimütigkeit haben zum Eintritt in das Heiligtum,
den er uns eröffnet hat als einen neuen und lebendigen Weg durch den Vorhang – das ist durch sein Fleisch –,
und einen großen Priester über das Haus Gottes,
so lasst uns hinzutreten mit wahrhaftigem Herzen in voller Gewissheit des Glaubens,
die Herzen besprengt ⟨und damit gereinigt⟩ vom bösen Gewissen und den Leib gewaschen mit reinem Wasser.
	\attrib{\bibverse{Hebrews}{10}{19-22}, \Cite{elb}}
\end{quote}
Jesu Opfer hat nicht nur den neuen Bund eingeleitet, Sünden vergeben und uns rein gemacht - er eröffnet uns auch den Zugang durch den Vorhang, in das Allerheiligste.
Mit unserem Verständnis des levitischen Systems ist auch diese Aussage direkt aus dem vorher gesagten logisch abzuleiten:
Wir sind rein gewaschen durch Jesu perfektes reinigendes Blut,
wir sind besprengt mit dem Blut des Bundes - und damit metaphysisch transformiert hin zu einem heiligen Volk Gottes.
So sind wir als Menschen also vorbereitet, das Heiligtum zu betreten.
Nicht nur wir, sondern auch das himmlische Heiligtum wurde vorbereitet durch Jesu Blut als Säuberungsopfer, mit dem das Heiligtum gereinigt wurde.
Alles ist bereit, alles ist in kultischer Ordnung und jetzt ist der Weg frei;
durch Jesu Leben, durch sein Blut, mit ihm als unserem Fürsprecher vollständig mit Gott versöhnt und ohne Angst in seiner Gegenwart.

Es folgt eine Ermahnung, an dem Bekentniss dieses Christus festzuhalten, gefolgt von dieser weiteren Warnung, dass Jesu endgültiges Opfer nicht verachtet werden darf:
\begin{quote}
Denn wenn wir mutwillig sündigen, nachdem wir die Erkenntnis der Wahrheit empfangen haben, bleibt kein Schlachtopfer für Sünden mehr übrig,
sondern ein furchtbares Erwarten des Gerichts und der Eifer eines Feuers, das die Widersacher verzehren wird.
Hat jemand das Gesetz Moses verworfen, stirbt er ohne Barmherzigkeit auf zwei oder drei Zeugen hin.
Wie viel schlimmere Strafe, meint ihr, wird der verdienen, der den Sohn Gottes mit Füßen getreten und das Blut des Bundes,
durch das er geheiligt wurde, für gemein erachtet und den Geist der Gnade geschmäht hat?
	\attrib{\bibverse{Hebrews}{10}{26-29}, \Cite{elb}}
\end{quote}
Hier zeigt sich wieder, was das Zentrum Jesu Wirkens ist: der neue Bund.
Denn seinem Bundesvolk in Christus hat der Vater völlig vergeben, sie angenommen und geheiligt.
Es gibt keine andere Möglichkeit, dem kommenden Gericht zu entkommen, als Teil dieses Bundes zu sein -
kein anderes Opfer ist in der Lage das zu ersetzen, was Christus auf viel bessere Art und Weise bewirkt hat.


Der Brief an die Hebräer ist für unser Verständnis des Opfers am Kreuz eine wahre Fundgrube an komplizierten Bezügen auf das levitische System und die prophetische Hoffnung
des neuen Bundes -
Bezüge, die plötzlich ein wunderschönes und klares Bild ergeben,
nachdem wir in den ersten Kapiteln die Farben und die Pinsel einzeln under die Lupe genommen haben.
Das Kreuz - das freiwillige Opfer des inkarnierten Wortes zur Vergebung unserer Sünden - ist ein Höhepunkt der Geschichte eines liebenden Gottes,
der seinem Bundesvolk großzügig und unverdient vergibt; der sein göttliches Blut des perfekten Lebens vergießt um uns rein zu waschen und zu heilen aus unserer korrupten Natur;
der seinen Sohn in den Tod gibt, damit sein Blut das himmlische Heiligtum reinigen kann, sodass wir unter der unaufhörlichen Fürsprache des immer lebenden Sohnes,
gewaschen im Heiligen Geist, und ohne Furch vor den Vater treten können.

\chapter{Gerechtigkeit und das letzte Gericht}
Ein wichtiger Bestandteil der klassischen protestantischen (und insbesondere reformierten) Soteriologie ist eine bestimmte
Ansicht darüber, auf welcher Basis die Gläubigen durch das letzte Gericht hindurch bestehen werden.
Die Kurzfassung ist diese:
Am Ende der Zeit wird jede Person vor Gottes Gericht erscheinen. Die nicht durch Jesus erlösten werden wegen ihrer Sünden zum Tode verurteilt.
Für die Erlösten sieht Gott statt ihrer eigenen Gerechtigkeit auf die Perfektion Jesu, die uns angezogen wurde und verdammt uns auf dieser Grundlage nicht.

Zu dieser Ansicht kommt das Verständnis hinzu, dass Gott jede Sünde bestrafen muss, und um uns nicht zu vernichtet stattdessen Jesus die Strafe auf sich nimmt
- diese zweite Seite der Medaillie ist dann PSA.
In diesem Artikel ist unsere eigentliche Frage die zweite: hat der Vater den Sohn für unsere Schuld bestraft?
Wenn wir diese Frage mit Nein beantworten, wie ich es tue, ist es aber nötig ein weiteres loses Ende zu erklären:
Was bedeutet es dann, dass wir durch Jesu Tod gerechtfertigt werden?
Dieser Frage wenden wir uns in diesem Kapitel vor. Gleichzeitig können wir beim beantworten dieser Frage
natürlich über einige Bibelstellen sprechen, die gerne in Verteidigungen von PSA verwendet werden - aber eine bessere Interpretation haben.

Hier ist meine Sicht vorweggenommen:
(a) Jeder Mensch wird im letzten Gericht für seine tatsächliche Gerechtigkeit nach seinen Werken gerichtet.
Die Gerechten werden leben, die Ungerechten werden vernichtet.
(b) Die griechischen Wörter hinter `Gerechtfertigung' und seiner verbalen und adjektivischen Form sind nicht ausschließlich forensisch, sondern
können ontologisch fungieren\footnote{
	Das heißt, `Gerechtigkeit', `gerecht machen', ... sind in manchen Kontexten valide Übersetzungen.
}
(c) Wenn Paulus diese Wörter im Römerbrief verwendet, sind sie besser in diesem ontologischen Charakter zu verstehen:
durch Glauben an Jesus werden wir gerecht gemacht, nicht bloß unabhängig von unseren Werken gerecht gesprochen.
Diese Konzepte so zu verstehen, und insbesondere den Römerbrief so zu lesen, ist die beste Möglichkeit, zwei Wahrheiten zu verstehen, die
häufig gegensätzlich erscheinen: Wir sind gerichtet nach Werken, aber gerechtfertigt durch Glauben.\footnote{
	Ich habe nicht genug Platz, um alle Verse zu besprechen, die in dieser Diskussion interessant wären.
	Ein gutes Beispiel ist \bibverse{Ephesians}{2}{8-9} (beachte zu diesem speziellen Vers: was meint Paulus hier mit `gerettet'?).
	Da mein primäres Ziel in diesem Kapitel ist, Paulus' Verständnis von Gerechtfertigung in Relation zu Jesu Tod und Leben zu verstehen
	bin ich hier nicht erschöpfend.
}

\section{\label{sec:gericht-nach-werken}Gericht nach Werken}
Die jüdische Literatur in der Zeit des zweiten Tempels ist in vielen speziellen Fragen recht divers.
Doch zumindest in dieser Frage gibt es keine Ausreißer: jeder Mensch wird am Ende nach seinen Werken gerichtet.
Wer gerecht gelebt hat wird leben, und wer ungerecht war wird vernichtet.
Wenn eine Person vor diesem Gericht besteht, dann weil sie tatsächlich gut gelebt hat.
In diesem Gericht sieht Gott die Person nicht an, und wird niemanden belohnen, der diese Belohnung nicht verdient hat.
Wie absolut und ausnahmslos dieses Verständnis ist hat VanLandingham in seinem Buch \Cite{judgement-and-justification-in-early-judaism-and-the-apostle-paul} gezeigt.
Hier gebe ich nur eine kurze Auswahl dieser Stellen wieder:

\begin{quote}
	Zerrissen wird die Erde,
und alles auf der Erde wird vergehen,
und alles wird gerichtet werden.
Doch mit den Frommen schließt Er Frieden
und schützt die Auserwählten
und Gnade waltet über ihnen.
Sie werden alle Gottes Eigentum
und sind im Glücke und gesegnet.
Er selber unterstützt sie alle
und Gottes Licht wird ihnen scheinen.
Er selber schließt mit ihnen Frieden.
Fürwahr! Er kommt mit Tausenden von Heiligen,
um über alle das Gericht zu halten
und alle Übeltäter zu vernichten
und alles Fleisch zurechtzuweisen
der schlimmen Taten wegen,
die sie so frevlerisch begingen,
sowie der kühnen Worte halber,
die gegen Ihn die Sünder frevelhaft gesprochen.\footnote{
	In \Cite[5]{aeth-hen} lesen wir weiterhin, dass die Gottes Gesetz nicht beachten `keinen Frieden finden' (4) werden, während für die Gerechten
	während oder nach dem Gericht gilt: `Und Weisheit wird verliehen den Auserwählten; sie leben all und sündigen nicht mehr'.(8)
	Wir werden später sehen, wie ähnlich diese Idee Paulus Argument im Römerbrief kommt, bloß dass Paulus das Ausgießen von Gottes Gerechtigkeit auf den Christus
	bezieht, nicht auf die Zeit nach dem letzten Gericht.
}
\attrib{\Cite[1:7-9]{aeth-hen}}
\end{quote}

\begin{quote}
Reflect on all these matters and seek from him that he may support your
counsel and keep far from you the evil schemings and the counsel of Belial,
so that at the end of time, you may rejoice in finding that some of our words
are true. And it shall be reckoned to you as justice when you do what is
upright and good before him, for your good and that of Israel.\footnote{
	Für weiteres Material aus Qumran, beachte z.B. die Zwei-Wege Kapitel in 1QS 3-4.
}
\attrib{\Cite[4QMMTe Frag. 14-17 Col II]{tdss-se}}
\end{quote}

\begin{quote}
	Wer rechtschaffen handelt, erwirbt sich Leben beim Herrn,
  und wer Unrecht thut, verwirkt selbst sein Leben in Verderben
	denn des Herrn Gerichte sind gerecht gegen Person und Haus.\footnote{
		Vergleiche auch \Cite[9:11]{pss-sol} (Des Herrn ist das Erbarmen über das Haus Israel immer und ewig.) mit \Cite[10:8]{pss-sol} (Des Herrn [Werk] ist die Erlösung über das Haus Israel, zur ewigen Freude.).
		Wie VanLandingham beschreibt ist hier die Gottes Erbarmen (seine Gnade) gleichbedeutend mit der Errettung als solche, nicht mit einer bestimmten Art und Weise auf die 
		die Erlösten anders als die Ungerechten im letzten Gericht behandelt werden. Vergleiche \bibverse{Isaiah}{56}{1} LXX für einen ähnlichen Parallelismus.
	}
	\attrib{\Cite[9:5]{pss-sol}}
\end{quote}

\begin{quote}
	Ich antwortete und sprach: Herr Gott, du hast ja in deinem Gesetze bestimmt,
	nur die Gerechten würden dies Erbteil bekommen, aber die Gottlosen sollten ins Verderben gehen.
	\attrib{\Cite[7:17]{fourth-ezra}}
\end{quote}

\begin{quote}
	Denn sieh, es kommen Tage;
da öffnet man die Bücher,
worin die Sünden aller Missetäter aufgeschrieben,
sowie die Kammern,
wo die Gerechtigkeit all derer aufgespeichert ist,
die in der Schöpfung recht gehandelt.
\attrib{\Cite[24:1]{second-baruch}}
\end{quote}

Viele dieser Texte basieren auf der Deuteronomischen Formel\bibfoot{Deuteronomy}{28}{}:
Wenn Israel Gott gehorcht werden sie gesegnet sein und erhaben sein - hier sehr praktisch als Nation auf dieser Welt.
Wenn sie aber nicht gehorchen, wird die Nation verflucht sein und ins Exil gehen.
Genau das geschieht, und Israel geht ins Exil - nicht weil Gott zu schwach ist, sondern weil das Volk gesündigt hat und Gott Gericht hält.
Selbst als Juda in sein Land zurückgekehrt ist sündigt das Volk weiter, und die Flüche des Bundes lasten weiter auf ihm.\footnote{
	\bibverse{Haggai}{1}{}, das Buch Maleachi, \bibverse{Zechariah}{11}{4-17}.
	In vielen anderen Stellen ringt Gott durch die Propheten um sein Volk, dass es nicht in die Wege ihrer Väter zurück fällt: \bibverse{Zechariah}{1}{1-6}.
	Und die Rückkehrer sahen ihre schwere Lage zumindest als Resultat der Sünde ihrer Väter: \bibverse{Nehemiah}{9}{32-37}, \bibverse{Ezra}{9}{5ff}.
	Wir haben in Kapitel \ref{chap:bundesfluch} bereits gesehen, dass die Juden zur Zeit des Neuen Testaments immer noch der Meinung waren, dass
	das Exil nicht völlig beendet ist.
}.

In dieser Situation sehen wir in Maleachi die Hoffnung, dass Gott an seinem Tag die gerechten und die gottlosen voneinander trennen wird\footnote{
	Vergleiche auch \bibverse{Isaiah}{65}{}, wo die Endzeitliche Hoffnung auf die Gerechten beschränkt ist und die Ungerechten nicht teilhaben.
	Ebenso \bibverse{Daniel}{12}{2-3}, in dem aber nur eine der Seiten (die Gerechten, die viele zur Gerechtigkeit führen) explizit beschrieben wird.
}:
\begin{quote}
	Denn siehe, der Tag kommt, der wie ein Ofen brennt. Da werden alle Frechen und alle, die gottlos handeln, Strohstoppeln sein. Und der kommende Tag wird sie verbrennen, spricht der HERR der Heerscharen, sodass er ihnen weder Wurzel noch Zweig übrig lässt.
	Aber euch, die ihr meinen Namen fürchtet, wird die Sonne der Gerechtigkeit aufgehen, und Heilung ist unter ihren Flügeln. Und ihr werdet hinausgehen und umherspringen wie Mastkälber.
	\attrib{\bibverse{Malachi}{3}{19-20}}
\end{quote}

Und auch im neuen Testament findet sich dieselbe Ansicht ohne große Diskussion\footnote{
	Und das sollten wir auch erwarten: Entweder der Jesus-Glauben reiht sich in das Verständnis aller anderen Juden der Zeit ein,
	oder wir erwarten eine klare Diskussion darüber, warum dieses Gericht doch anders aussehen wird.
}:
\begin{quote}
	Wenn aber der Sohn des Menschen kommen wird in seiner Herrlichkeit und alle Engel mit ihm, dann wird er auf seinem Thron der Herrlichkeit sitzen;
	und vor ihm werden versammelt werden alle Nationen, und er wird sie voneinander scheiden, wie der Hirte die Schafe von den Böcken scheidet.
	Und er wird die Schafe zu seiner Rechten stellen, die Böcke aber zur Linken.
	Dann wird der König zu denen zu seiner Rechten sagen: Kommt her, Gesegnete meines Vaters, erbt das Reich, das euch bereitet ist von Grundlegung der Welt an!
	Denn mich hungerte, und ihr gabt mir zu essen; mich dürstete, und ihr gabt mir zu trinken; ich war Fremdling, und ihr nahmt mich auf;
	nackt, und ihr bekleidetet mich; ich war krank, und ihr besuchtet mich; ich war im Gefängnis, und ihr kamt zu mir.
	Dann werden die Gerechten ihm antworten und sagen: Herr, wann sahen wir dich hungrig und speisten dich? Oder durstig und gaben dir zu trinken?
	Wann aber sahen wir dich als Fremdling und nahmen dich auf? Oder nackt und bekleideten dich?
	Wann aber sahen wir dich krank oder im Gefängnis und kamen zu dir?
	Und der König wird antworten und zu ihnen sagen: Wahrlich, ich sage euch, was ihr einem dieser meiner geringsten Brüder getan habt, habt ihr mir getan.
	Dann wird er auch zu denen zur Linken sagen: Geht von mir, Verfluchte, in das ewige Feuer, das bereitet ist dem Teufel und seinen Engeln!
	Denn mich hungerte, und ihr gabt mir nicht zu essen; mich dürstete, und ihr gabt mir nicht zu trinken;
	ich war Fremdling, und ihr nahmt mich nicht auf; nackt, und ihr bekleidetet mich nicht; krank und im Gefängnis, und ihr besuchtet mich nicht.
	Dann werden auch sie antworten und sagen: Herr, wann sahen wir dich hungrig oder durstig oder als Fremdling oder nackt oder krank oder im Gefängnis und haben dir nicht gedient?
	Dann wird er ihnen antworten und sagen: Wahrlich, ich sage euch, was ihr einem dieser Geringsten nicht getan habt, habt ihr auch mir nicht getan.
	Und diese werden hingehen zur ewigen Strafe, die Gerechten aber in das ewige Leben.
	\attrib{\bibverse{Matthew}{25}{31-46}, \Cite{elb}}
\end{quote}
Diese Gerichtszene auf Jesu Lippen trägt alle üblichen Charakterisika von Szenen über das letzte Gericht:
Gott - hier der Menschensohn - hält Gericht mit seinen Engeln. Die Menschen werden in zwei Gruppen eingeteilt: die Gerechten und die Ungerechten.
Diese Aufteilung der Menschen geschieht basierend auf Werken - hier speziell basierend auf der Nächstenliebe, die jeder gezeigt hat.
Die einzige Neuerung in dieser Beschreibung ist, dass der Menschensohn das Gericht hält, nicht der Vater selbst.

\begin{quote}
	Wundert euch darüber nicht, denn es kommt die Stunde, in der alle, die in den Gräbern sind, seine Stimme hören
	und hervorkommen werden; die das Gute getan haben zur Auferstehung des Lebens,
	die aber das Böse verübt haben zur Auferstehung des Gerichts.
	\attrib{\bibverse{John}{5}{28-29}, \Cite{elb}}
\end{quote}
Wieder gibt es ein Gericht über alle Menschen, die Entscheidung fällt basierend auf dem Guten oder Bösen, dass jeder getan hat, und das Ergebnis ist Leben oder Gericht.

\begin{quote}
	Denn wir müssen alle vor dem Richterstuhl Christi offenbar werden,
	damit jeder empfängt, was er durch den Leib \textlangle vollbracht\textrangle,
	dementsprechend, was er getan hat, es sei Gutes oder Böses.
	\attrib{\bibverse{IICorinthians}{5}{10}, \Cite{elb}}
\end{quote}
Hier sehen wir, dass Paulus sich wie alle anderen Juden seiner Zeit in diesen Konsens einreiht:
Es gibt ein endgültiges Gericht, dort wird jeder für seine Taten
gerichtet werden. Die einzige Inovation ist, dass Christus der Richter ist, nicht der Vater.
Am deutlichsten finden wir diese Wahrheit in Paulus Briefen im zweiten Kapitel des Römberbriefs.
Auch in dieser Stelle weicht Paulus nicht von der gewöhnlichen Ansicht der Juden zum ewigen Gericht ab: 
Gericht am Ende der Zeit, basierend auf Werken, mit dem Resultat von Leben oder Gottes endgültige Strafe.

\begin{quote}
	Nach deiner Störrigkeit und deinem unbußfertigen Herzen aber häufst du dir selbst Zorn auf für den Tag des Zorns und der Offenbarung des gerechten Gerichts Gottes,
	der einem jeden vergelten wird nach seinen Werken:
	denen, die mit Ausdauer in gutem Werk Herrlichkeit und Ehre und Unvergänglichkeit suchen, ewiges Leben;
	denen jedoch, die von Selbstsucht \textlangle bestimmt\textrangle und der Wahrheit ungehorsam sind, der Ungerechtigkeit aber gehorsam, Zorn und Grimm.
	Bedrängnis und Angst über die Seele jedes Menschen, der das Böse vollbringt, sowohl des Juden zuerst als auch des Griechen;
	Herrlichkeit aber und Ehre und Frieden jedem, der das Gute wirkt, sowohl dem Juden zuerst als auch dem Griechen.
	\attrib{\bibverse{Romans}{2}{5-10}, \Cite{elb}}
\end{quote}

Wir können also zusammenfassend sagen:
Beim letzten Gericht wird jeder Mensch nach seinen Taten gerichtet. Die Gerechten werden leben und die, die Gott ablehnen werden vernichtet.
Über diese Tatsache gibt es in der jüdischen Literatur bis ins erste Jahrhundert keine Diskussion, nicht einmal im Brief des Paulus an die Römer.
Wie passt das Gericht nach Werken und die `Gerechtfertigung' durch Glauben zusammen?\footnote{
	Innerhalb von 2 Kapiteln in Römer 2-3!
}
Um diese Frage beantworten zu können wenden wir uns jetzt den griechischen Wörtern von der Wurzel δικαι zu: δικαιος, δικαιοσυνη, und δικαιοω.
Wir werden sehen, dass sich der scheinbare Widerspruch nicht auflöst, indem wir das Gericht nach Werken umdefinieren, sondern indem
wir die δικαι-Wörter, und damit das Konzept der `Gerechtfertigung' im Neuen Testament besser verstehen.

\section{Die δικαι-Wörter}
Die δικαι-Wörter, insbesondere in den paulinischen Briefen, sind in der aktuellen wissenschaftlichen Literatur äußerst
umstritten. Bei der Frage `was meint Paulus mit Gerechtfertigung' existieren mindestens die folgenden Ansicht:
`Gerechtsprechung'\footnote{
	Die klassische protestantische Interpretation. Siehe zum Beispiel \Cite{both-judge-and-justifier}, \Cite[82ff]{nicnt-romans}.
}, `Deklaration als Mitglied des Bundes'\footnote{
	Verbunden mit der New Perspective on Paul, insbesondere gehalten von NT Wright.
}, `Wiederherstellung der Bundesbeziehung, verbunden mit der Hoffnung auf Freispruch'\footnote{
	\Cite{inhabiting-the-cruciform-god}. Diese Sicht ist mit NT Wrights Sicht verbunden - aber Wright schließt einen transformativen
	Teil in der Bedeutung von δικαιοω explizit aus (\Cite{paul-and-the-faithfulness-of-god}), während Gorman ihn mit einbezieht.
}.
Anderswo ist `justification inseparable from glorification'\footnote{
	\Cite[197]{paul-a-new-covenant-jew}. Pitre, Barber, und Kincaid sind explizit der Meinung, dass Gerechtfertigung ausschließlich als legale
	Transaktion nicht ausreichend ist.
}.

Bei der Betrachtung jedes Wortes ist es fundamental wichtig, zwischen der invarianten Bedeutung - also der Bedeutung, die ein Wort in jeden Kontext einbringt - und
der Interpretation des Wortes in einem Kontext - also wie die invariante Bedeutung mit dem Kontext zusammen interpretiert wird - zu unterscheiden.
Natürlich reichen diese Seiten nicht aus, um mit allen Positionen zu den δικαι-Wörtern im Detail zu interagieren.
Stattdessen möchte ich zeigen, dass
(a) die δικαι Wortgruppe in ihrer invarianten Bedeutung nicht auf eines der Extreme Ontologie oder Forensik beschränkt ist, und
(b) dass eine genaue Untersuchung des Kontexts nötig ist, um zwischen verschiedenen Interpretationen zu unterscheiden.
Ich werde auch eine Hypothese dafür geben, wie diese Disambiguierung im allgemeinen vorgenommen werden kann.

Das linguistisch komplexeste der δικαι-Wörter ist die verbale Form. Im Gegensatz dazu sind das adjektiv δικαιος und das deadjektivische Substantiv δικαιοσυνη deutlich einfacher zu analysieren:
\begin{exe}
	\ex \label{ex:dikaios-open}
	{
		ὁ δὲ ἄνθρωπος, ὃς ἔσται δίκαιος, ὁ ποιῶν κρίμα καὶ δικαιοσύνην, [...], δίκαιος οὗτός ἐστιν \\
		Und der Mensch, der gerecht ist, der Recht und Gerechtigkeit ausübt, [es folgt eine Liste von bösen Dingen, die dieser Mensch nicht tut], gerecht ist er.
	}
	\attrib{\bibverse{Ezekiel}{18}{5-9} LXX}
\end{exe}
An diesem Beispiel sehen wir gut, wie das Adjektiv und das Substantiv parallel verwendet werden und beide einen Menschen charakterisieren, der richtig handelt.

\begin{exe}
	\ex \label{ex:dikaios-gradable}
		{Δίκαιος σὺ ὑπὲρ ἐμέ \\
		 `du warst gerechter als ich'
		}
	\attrib{\bibverse{ISamuel}{24}{18} = LXX 1. Buch der Könige 24:18}
\end{exe}
Hier wird sichtbar, dass δικαιος für Vergleiche verwendet werden kann.
Adjektive, die für Vergleiche verwendet werden können nennt man auch \textit{gradierbar}: sie können zu bestimmten Graden zutreffen.
Beispiele im deutschen sind `groß' und `feucht', aber zum Beispiel nicht `tot'.
Gradierbare Adjektive haben einige besondere linguistische Eigenschaften, insbesondere hat die Art der \textit{Skala}, die sie beschreiben einen großen Einfluss auf ihr
syntaktisches Verhalten und ihre Interpretation.\footnote{
	Für eine genauere Einführung, siehe \Cite{scale-structure}.
}
Vergleichen wir beispielsweise diese deutschen Beispiele:
\begin{exe}
\ex\label{ex:open-closed} 
\begin{xlist}
	\ex[]{Max ist nicht groß, aber größer als Moritz.}
	\ex[]{\label{ex:feucht-gradable}Nach dem Schwimmen ist meine Badehose feuchter als mein Handtuch.}
	\ex[\texttt{\#}]{\label{ex:upper-half-open}Meine Hand ist nicht feucht, aber es ist etwas Wasser auf ihr.\footnote{\texttt{\#} bedeutet: Dieses Beispiel ist grammatisch erlaubt, aber in seiner Bedeutung widersprüchlich.}}
	\ex[*] {\label{ex:non-gradable}Bob ist nicht lebendig, aber Alice ist töter.\footnote{* bedeutet: Dieses Beispiel ist ungrammatisch.}}
\end{xlist}
\end{exe}
Wir sehen zunächst, dass manche adjektive nicht gesteigert und verglichen werden können: `tot' ist nicht gradierbar, und Beispiel \ref{ex:non-gradable} ist ungrammatisch.
Andere adjektive sind steigerbar und auf einer \textit{offenen Skala}. Das eine Person `groß' ist, entscheidet sich nur im Vergleich mit anderen Personen:
Sind Max und Moritz im selben Raum, so ist in diesem Kontext Max groß, und Moritz klein. Es gibt keine Untergrenze und Obergrenze für das `Großsein'.
Man sagt, dass die Skala beidseitig offen ist.
Beispiel \ref{ex:upper-half-open} zeigt, dass sich das deutsche adjektiv `feucht' anders verhält.
Dieses adjektiv ist gradierbar, wie \ref{ex:feucht-gradable} zeigt, aber es gibt eine fixe Untergrenze für `Feuchtheit'.
Sobald auch nur etwas Wasser auf einer Oberfläche ist, ist diese Oberfläche automatisch feucht.

\begin{exe}
\ex\label{ex:ambiguous-scale} 
\begin{xlist}
	\ex[]{\label{ex:dry-ambiguous}Meine Hände sind trocken.}
	\ex[]{\label{ex:dry-closed} Ich habe meine Hände gerade abgetrocknet und sie sind trocken.}
	\ex[]{\label{ex:dry-open} Meistens sind meine Hände trocken, aber unter Stress schwitzen sie schnell.}
\end{xlist}
\end{exe}

Beispiel \ref{ex:ambiguous-scale} schließlich zeigt, dass ein und dasselbe Wort `trocken' je nach Kontext verschiedene Skalenstrukturen realisieren kann:
In \ref{ex:dry-closed} ist die Skala geschlossen: eine individuelle Situation ist im Blick, und `meine Hände sind trocken' bedeutet
`es befindet sich aktuell keine Flüssigkeit auf ihnen' - das Maximum an Trockenheit ist in diesem Moment tatsächlich erreicht.
In \ref{ex:dry-open} wird eine allgemeine Qualität der Hände beschrieben.
Verglichen mit den Händen anderer Menschen sind meine Hände üblicherweise trocken.
Trotzdem können meine Hände in speziellen Situationen naß sein. Die Skala ist also offen, es wird nicht über ein Maximum an Trockenheit geredet.
Zuletzt ist die Aussage in \ref{ex:dry-ambiguous} unklar. Das deutsche Wort `trocken' hat weder eine offene Skala (allgemeine Qualität), noch eine geschlossene Skala
(Zustand in einer bestimmten Situation) als Teil seiner invarianten Bedeutung.

Mit dieser Theorie kommen wir nun zurück zu den δικαι-Wörtern.
In \ref{ex:dikaios-gradable} sahen wir bereits, dass δικαιος gradierbar ist.
Die gesamte Wortgruppe markiert ihre Skalenstruktur nicht, analog zum deutschen Wort `trocken'.\footnote{
	Hier liegt eines der Probleme, in deutscher oder englischer Sprache über die `Gerechtfertigung' zu reden:
	Im Deutschen und Englischen haben wir kein Verb, das in seiner Skalenstruktur so flexibel ist wie das griechische δικαιοω.
	Wir müssen also beim Übersetzen eine Entscheidung darüber treffen, ob das griechische Wort geschlossen oder offen zu interpretieren ist.
}
Das Adjektiv δικαιος kann eine offene Skala haben, und die grundsätzliche Art eines Menschen beschreiben, wie in \ref{ex:dikaios-open} oder in folgendem Beispiel:
\begin{exe}
	\ex
		{
			ἦσαν δὲ δίκαιοι ἀμφότεροι ἐναντίον τοῦ θεοῦ, πορευόμενοι ἐν πάσαις ταῖς ἐντολαῖς καὶ δικαιώμασιν τοῦ κυρίου ἄμεμπτοι. \\
			Beide aber waren gerecht vor Gott und wandelten untadelig in allen Geboten und Satzungen des Herrn.
		}
		\attrib{\bibverse{Luke}{1}{6}\footnote{
			Es ist wahr, dass die `Gebote und Satzungen des Herrn' einen Standard bilden, der die Skala nach oben abschließen könnte.
			In dieser Stelle im Diskurs wird allerdings der Charakter von Zacharias und Elisabeth beschrieben,
			keine individuelle Situation, in der ihr Handeln in Frage gestellt wird.
			Stattdessen ist ihr `Wandeln in den Geboten' ebenso eine allgemeine Charaktereigenschaft - aus theologischen Gründen wollen
			wir diese Verse so oder so nicht verstehen als wären Zacharias und Elisabeth sündlos - also in jeder individuellen Sitution ihres Lebens gerecht.
	}}
\end{exe}
Ein Beispiel, in dem die Skala geschlossen ist findet sich in \bibverse{IKings}{8}{31-32} (=LXX 3. Buch der Könige 8:31-32),
wo eine individuelle Sünde von einem Menschen gegen einen anderen in Sicht ist, und der Unschuldige δικαιος genannt wird.
Es ist nicht der Charakter der Personen, sondern eine individuelle Tat in Frage.\footnote{
	δικαιος wird sehr selten in dieser geschlossenen Form verwendet, viel häufiger taucht es mit einer offenen Skala auf.
	Der Grund für diese Tatsache ist wahrscheinlich, dass αθωος (unschuldig) zur Verfügung steht, das für eine geschlossene Skala markiert ist.
	Wir erwarten, und sehen in den Daten, dass das unmarkierte δικαιος nur selten in einer Art und Weise verwendet wird, die vom markierten αθωος besser gefüllt wird.
	Ein gutes minimales Beispiel bildet \bibverse{Matthew}{27}{24}.
	Pilatus ist unschuldig an einer individuellen Tat (geschlossene Skala), sein allgemeiner Charakter (offene Skala) ist hier nicht im Spiel und αθωος wird verwendet.
	Die verbale Form αθωοω ist deutlich seltener zu finden (nicht im NT, nur LXX), und damit korreliert wahrscheinlich die Tatsache, dass δικαιοω häufiger für die
	geschlossene Skala verwendet wird als sein Adjektiv. Siehe für ein klares Beispiel \ref{ex:dikaiow-closed}.
}
Das Substantiv δικαιοσυνη unterscheidet sich kaum vom Adjektiv.
Wenn wir nun zur verbalen Form δικαιοω kommen ensteht ein ähnliches Bild, aber die offene Skala wird häufiger verwendet als beim Adjektiv:
\begin{exe}
	\ex\label{ex:dikaiow-closed}
	\begin{xlist}
		\ex \label{ex:dikaiow-not-transformational} {
			Ἐὰν δὲ γένηται ἀντιλογία ἀνὰ μέσον ἀνθρώπων [...] καὶ δικαιώσωσιν τὸν δίκαιον καὶ καταγνῶσιν τοῦ ἀσεβοῦς \\
			Wenn es zu einem Rechtsstreit zwischen Männern kommt [...], und sie sprechen den Gerechten gerecht, und den Schuldigen für unrecht befinden
		}
		\attrib{\bibverse{Deuteronomy}{25}{1} LXX}
		\ex \label{ex:dikaiow-forensic-real} {
			ἐκ γὰρ τῶν λόγων σου δικαιωθήσῃ καὶ ἐκ τῶν λόγων σου καταδικασθήσῃ \\
			denn aus deinen Worten wirst du gerecht befunden werden, und aus deinen Worten wirst du verdammt werden
		}
		\attrib{\bibverse{Matthew}{12}{37}\footnote{
				Es ist hier (a) das Gericht im Blick, (b) stehen individuelle Worte im Fokus, anhand derer die Menschen gerichtet werden.
		}}
	\end{xlist}
	\ex\label{ex:dikaiow-open}
	\begin{xlist}
		\ex {
				Ἄρα ματαίως ἐδικαίωσα τὴν καρδίαν μου καὶ ἐνιψάμην ἐν ἀθῴοις τὰς χεῖράς μου \\
				Jetzt habe ich umsonst mein Herz gerecht gemacht, und in Unschuld meine Hände gewaschen!
		}
		\attrib{\bibverse{Psalms}{73}{13} = LXX Psalms 72:13}
		\ex \label{ex:dikaiow-transformational} {
			καὶ ταῦτά τινες ἦτε ἀλλʼ ἀπελούσασθε, ἀλλʼ ἡγιάσθητε, ἀλλʼ ἐδικαιώθητε \\
			Und das [Sünder] sind einige von euch gewesen, aber ihr seid abgewaschen, aber ihr seid heilig gemacht, aber ihr seid gerecht gemacht
		}
		\attrib{\bibverse{ICorinthians}{6}{11}}
	\end{xlist}
\end{exe}
In den Beispielen \ref{ex:dikaiow-closed} ist die Skala geschlossen, und das Verb muss dementsprechend interpretiert werden:
\ref{ex:dikaiow-not-transformational} erlaubt keine Interpretation außer,
dass die Person in Frage in diesem individuellen Rechtstreit Recht erhält oder für gerecht befunden wird.
In beiden Beispielen in \ref{ex:dikaiow-open} muss der allgemeine Charakter oder der generelle Zustand der Person (bzw. des Herzens) in Frage sein,
da keine individuelle Tat oder ein spezieller Streit das Thema ist.
Gerade in \ref{ex:dikaiow-transformational} ist keine Interpretation möglich außer, dass δικαιοω hier parallel zu `abgewaschen' und `heilig gemacht' steht:
für eine Transformation des allgemeinen Charakters der Person.
Die Interpretationen in \ref{ex:dikaios-open} sind notwendigerweise inchoativ (sie markieren den Beginn eines Zustands).
Die Interpretationen mit geschlossener Skala in \ref{ex:dikaiow-closed} dagegen beschreiben das ausdrücken oder feststellen einer
forensischen Realität durch das Aufrichten oder für gerecht Erklären des Gerechten.\footnote{
	Vergleiche auch mit \Cite{hebrew-verbs-of-judgement}.
	Die Situation ist bei den hebräischen und griechischen Wörtern aus dieser Domäne sehr ähnlich.
}
Zusammenfassend können wir also sagen, dass die δικαι-Wörter unmarkiert für die Skalenstruktur sind.
Der Kontext muss disambiguieren, ob
(a) eine geschlossene Skala, also eine individuelle Tat, und damit die Schuld oder Unschuld oder die Erfüllung eines fixen Rechtsstandards, oder
(b) eine offene Skala, also der Charakter oder die allgemeine Disposition einer Person beschrieben wird.

Um es theologischer auszudrücken:
Eine Theorie, die eine fixe theologische Bedeutung für die δικαι-Wörter fixieren möchte ist zum Scheitern verurteilt.
In jedem Fall ist der Kontext der individuellen Stelle nötig, um zwischen grundlegend verschiedenen Interpretationen unterscheiden zu können.
Werden den Geretteten ihre Sünden nicht angerechnet? Ja.
Formen die Geretteten im Christus den neuen Bund, dessen Mitglieder in rechter Beziehung zu Gott stehen? Ja.
Werden die Gläubigen aus ihrem Leben der Sünde heraus gerissen, und vom Geist befähigt Gott wohlgefällig zu leben? Ja.
Aber keine dieser Aussagen darf mit dem Wort `Gerechtfertigung' gleichgesetzt werden. Jeder Text, und selbst jede Phrase in einem Text muss für sich stehen und gelesen werden
um zu einer belastbaren Theologie zu gelangen.

\section{Gerechtigkeit und das Kreuz im Brief an die Römer}
In seinem Brief an die Römer beschreibt der Apostel Paulus so ausführlich wie kaum anderswo wie Gott die Menschen
durch das Wirken seines Sohne aus ihrer miseren Lage holt.
Da wir hier in einem anhaltenden Argument über das letzte Gericht, Jesu Tod am Kreuz, unseren neuen Frieden mit Gott, und unser neues Leben in Christus hören
möchte ich im folgenden zeigen, wie das neue Verständnis der δικαι-Wörter hilft, den Römerbrief konsistent zu interpretieren,
und wie in diesem Verständnis Jesu Tod in stellvertretender Übernahme unserer Schuld keine Rolle spielt.

\subsection{Römer 2}
Mit \bibverse{Romans}{2}{} beginnen wir sofort in einer äußerst umstrittenen, aber gleichzeitig oft ignorierten Passage des Briefe.
In Abschnitt \ref{sec:gericht-nach-werken} sahen wir bereits, dass dieses Kapitel eine der Stellen im Neuen Testament ist, in der das endgültige Gericht
nach Werken beschrieben wird.
Die größte Interpretative Frage zu diesem Kapitel ist nicht, was es bedeutet, sondern ob es konstruiert ist und damit eine hypothetische Situation beschreibt:
ein Gericht nach Werken, vor dem die Gläubigen gerettet werden.
Ist es also  Paulus' Argument, dass vor diesem Gericht, wenn es wirklich nach Werken geschieht, kein Mensch gerecht ist?

Diese Sicht ist äußerst unplausibel. Sicherlich ist die Klasse der gerechten Menschen in Texten über das letzte Gericht
klein: die meisten Menschen, selbst die meisten Juden, sind ungerecht und werden im letzten Gericht vernichtet.
Doch die klassische Position, dass kein Mensch vor Gottes Gericht gerecht ist findet sich nirgends sonst\footnote{
	\bibverse{Luke}{1}{6} war bereits zitiert: hier sind zwei Menschen `gerecht und wandeln tadellos in allen Gesetzen und Satzungen des Herrn'.
	Hat sich Lukas versprochen?
}, und erzeugt auch keine konsistente Interpretation dieses Kapitels.
Um das zu zeigen, betrachten wir zunächst v13:
\begin{quote}
	es sind nämlich nicht die Hörer des Gesetzes gerecht vor Gott, sondern die Täter des Gesetztes werden gerecht gesprochen werden.
	\attrib{\bibverse{Romans}{2}{13}}
\end{quote}
Es ist in der ganzen Passage (a) das letzte Gericht, und (b) ein klarer Standard in Sicht: Gottes Wille im Gesetz offenbart.
Damit ist die Skalenstruktur des Verbs δικαιοω geschlossen und eine entsprechende Interpretation wie hier angegeben ist nötig.
Das Ausführen des Gesetzes (der Torah) ist nötig, um von Gott im Gericht als gerecht befunden zu werden.
\begin{quote}
	Denn wenn Nationen, die von Natur aus kein Gesetz haben\footnote{
		Diese Übersetzung ist mit großer Sicherheit korrekt, kontra \Cite{nicnt-romans}.
		Die griechische Phrase ist ἔθνη τὰ μὴ νόμον ἔχοντα φύσει τὰ τοῦ νόμου ποιῶσιν, und die Frage ist, ob φύσει die erste Phrase `Nationen, die das Gesetz nicht haben',
		oder die zweite Phrase `sie tun das Gesetz' modifiziert.
		(a) In adverbialer Verwendung modifiziert φύσει Zustände, keine Ereignisse (in NT und LXX: \bibverse{Galatians}{2}{15} mit Null-Copula, \bibverse{Galatians}{4}{8}, \bibverse{Ephesians}{2}{3}, \bibverse{Wisdom}{13}{1} mit Null-Copula).
		(b) Die Idee, dass Heiden von Natur aus dem Gesetz entsprechend handeln ist inkonsistent mit \bibverse{Galatians}{2}{15}. Heiden sind von Natur aus Sünder, und haben von Natur aus das Gesetz nicht.
	}, das Gesetz tun,
	so sind diese, die kein Gesetz haben, sich selbst ein Gesetz.
	Sie beweisen, dass das Werk des Gesetzes in ihren Herzen geschrieben ist
	\attrib{\bibverse{Romans}{2}{14-15}}
\end{quote}
Es gibt also Heiden, die zwar von Natur aus kein Gesetz haben, aber doch die Dinge des Gesetzes tun.
So beweisen sie, dass Gottes Gesetz auf ihren Herzen geschrieben ist - die prophetische Hoffnung, die wir schon in Kapitel \ref{chap:prophetische-hoffnung} gesehen haben
und die jetzt durch den Neuen Bund im Christus Jesus offenbart ist.
Und aus diesem Grund können diese Heiden die Juden mit denen Paulus sein Argument führt anklagen: Sie tun das Gesetzt, in Christus, und werden deshalb im Gericht bestehen.
Die Juden ruhen sich darauf aus, die Torah zu kennen aber erlangen doch Abseits von Christus keine Gerechtigkeit.
Diese Verse zeigen, dass es für Paulus sehr wohl möglich ist für einen Menschen, das Gesetz zu erfüllen - nämlich in Christus.
Genau zu diesem Punkt werden wir in der Diskussion über \bibverse{Romans}{8}{3} zurückkehren: ``die Rechtsforderung des Gesetzes wird erfüllt in uns, die wir nicht nach dem Fleisch, sondern nach dem Geist wandeln''.

Ein sehr ähnliches Argument finden wir am Ende des Kapitels:
Paulus beschuldigt Juden, die sich auf das Gesetz stützen und doch Böses tun (v17-24), und setzt im Kontrast dazu den Heidenchristen, beschnitten im Herz statt am Körper,
dem diese Beschneidung tatsächlich etwas nützt:
\begin{quote}
	Wenn nun der Unbeschnittene die Rechtsforderungen des Gesetzes befolgt, wird nicht sein Unbeschnittensein für Beschneidung gerechnet werden,
	und das Unbeschnittensein von Natur, das das Gesetz erfüllt, dich richten, der du mit Buchstaben und Beschneidung ein Gesetzesübertreter bist?
	\attrib{\bibverse{Romans}{2}{26-27}, \Cite{elb}}
\end{quote}
Diese Argument ist absolut sinnlos, wenn Paulus hier über eine hypothetische Person redet, in dem Wissen dass sie nicht existiert, und
dass er im nächsten Kapitel erklären wird, dass diese Person nicht existiert.
Paulus legt hier eine sehr reale Situation auf den Tisch: der Heide, der abseits der körperlichen Beschneidung die Forderung des Gesetzes erfüllt
klagt den untreuen Juden an, den Paulus zur Umkehr bewegen will.\footnote{
	Moos Argument `But Paul would not describe a Christian as having been granted membership among the saved as a result of obeying the law' \Cite{nicnt-romans}
	ist zirkulär, und verliert seine Kraft wenn wir im folgenden sehen werden, dass `Gerechtfertigung aus Glauben' nicht `im letzten Gericht unabhängig von Werken
	gerecht gesprochen' bedeutet.
	Wenn Paulus das letzte Gericht so versteht wie jeder andere jüdische Text seiner Zeit, und so wie seine eigene Beschreibung auf den ersten Blick aussieht,
	dann ist `Christen werden im letzten Gericht gerecht erklärt, weil sie wirklich gerecht waren' exakt Paulus' Meinung.
}
Die Tuer des Gesetzes werden in Gottes Gericht bestehen.
Es stellen sich jetzt zwei Fragen: erstens, wie kann es sein, dass die Juden die das Gesetz haben nicht dadurch Tuer des Gesetzes sind.
Zweitens, wie kann es sein, dass die Heiden ohne Gesetz Tuer des Gesetzes werden?\footnote{
	Für eine detailliertere Diskussion über verschiedene Meinungen zu Römer 2 und eine anhaltendere Verteidigung der realen Position
	siehe \Cite[215ff]{judgement-and-justification-in-early-judaism-and-the-apostle-paul}.
}

\subsection{Römer 3:9-20}
Die tatsächliche Grundlage dafür, dass viele Kommentatoren Kapitel 2 auf die eine oder andere Art und Weise als hypothetisch abschreiben ist,
dass sie einen Widerspruch zwischen Paulus' Aussagen in Kapitel 3 sehen. Dieser Stelle wenden wir uns nun zu.
\begin{quote}
	Was nun? Haben wir einen Vorzug? Durchaus nicht! Denn wir haben sowohl Juden als auch Griechen vorher beschuldigt, alle unter der Sünde zu sein,
	wie geschrieben steht: ``Da ist kein Gerechter, auch nicht einer; da ist keiner, der verständig ist; da ist keiner, der Gott sucht.
	Alle sind abgewichen, sie sind allesamt untauglich geworden; da ist keiner, der Gutes tut, da ist auch nicht einer.''
	[...]
	% ``Ihr Schlund ist ein offenes Grab; mit ihren Zungen handelten sie trügerisch.'' ``Viperngift ist unter ihren Lippen.''
	% ``Ihr Mund ist voll Fluchens und Bitterkeit.''
	% ``Ihre Füße sind schnell, Blut zu vergießen; Verwüstung und Elend ist auf ihren Wegen, und den Weg des Friedens haben sie nicht erkannt.''
	% ``Es ist keine Furcht Gottes vor ihren Augen.''
	Wir wissen aber, dass alles, was das Gesetz sagt, es denen sagt, die unter dem Gesetz sind,
	damit jeder Mund verstopft wird und die ganze Welt dem Gericht Gottes verfallen ist.
	Darum: Aus Gesetzeswerken wird kein Fleisch vor ihm gerechtfertigt werden;
	denn durchs Gesetz \textlangle kommt\textrangle\ Erkenntnis der Sünde.
	\attrib{\bibverse{Romans}{3}{9-20}, \Cite{elb}}
\end{quote}
Im Kontrast zu \bibverse{Romans}{2}{13} ``die Täter des Gesetzes werden gerecht gesprochen werden'' bildet v20 eine Interpretatives Rätsel.
Besonders die protestantische Position löst dieses Problem häufig auf, indem 2:13 hypothetisch gelesen wird: Im Prinzip würden die Täter des Gesetzes schon
gerecht gesprochen werden, aber Paulus zeigt mit seinen Zitaten aus dem Alten Testament, dass kein einziger Mensch diesen Standard tatsächlich erfüllt.
Da ich im letzten Abschnitt dafür plädiert habe, dass Römer 2 nicht hypothetisch uminterpretiert werden kann muss sich dieses Paradox also in 3:20 auflösen.
Die primäre Erklärung, mit der sich dieses Paradox auflöst ist die Feststellung, dass δικαιοω in v20 auf einer offenen Skala agiert, dass Paulus
also über den allgemeinen Zustand der Menschen, nicht die individuellen Sünden des Lebens am letzten Gericht spricht.
Entsprechend ist meine Interpretation, dass `aus Werken des Gesetzes kein Mensch gerecht wird'.
Für diese Sicht sprechen folgende Punkte:

Erstens: Die Bibelstellen, die Paulus für sein Argument zu Rate zieht sagen nicht aus, dass keine einzige Person vor Gottes Gericht bestehen kann.
Sie zeigen den natürlichen Hang aller Menschen - auch der Juden - zur Sünde.
\bibverse{Ecclesiastes}{7}{21} sagt aus, dass kein Mensch so gerecht ist, dass er nie sündigt.
Aber es ist in der jüdischen Literatur klar, dass `Gericht nach Werken' nicht `jeder der auch nur einmal etwas Böses getan hat wird untergehen' bedeutet.
Sündlose Perfektion ist nicht der Rahmen in dem Gericht nach Werken zu verstehen ist.\footnote{
	Siehe zum Beispiel \Cite[CD-A Col II = 4Q266 2 II]{tdss-se}: ``God loves knowledge; he has established wisdom and counsel
	before him; prudence and knowledge are at his service; patience is his and
	abundance of pardon, to atone for those who repent from sin; however,
	strength and power and a great anger with flames of fire by the \textlangle hand\textrangle\ of all
	the angels of destruction against those turning aside from the path and abominating the precept''
	Selbst in der sehr streng legalistischen Gemeinschaft in Qumran ist axiomatisch akzeptiert, dass Menschen umkehren und sich der Wahrheit zuwenden können.
}
Einige der anderen zitierten Stellen beschreiben, dass Menschen vor Gott gerecht oder unschuldig ist, und Menschen deswegen Gnade brauchen.\footnote{
	\bibverse{Psalms}{143}{2} (der Author), \bibverse{Isaiah}{59}{7-8} (Israel).
	Diese Aussage widerspricht meiner Interpretation nicht: selbstverständlich sündigt jeder Mensch, und natürlich ist es Gottes Gnade,
	die dem Gerechten seine Schuld vergibt.
	Mein Punkt ist: wir sollten solch negative Anthropologie nicht soweit strapazieren, dass jede Gerichtsszene
	zu einer Farce wird, in der niemand auf der Seite Gottes steht.
	Sieht Jesus in \bibverse{Matthew}{25}{31-46} oder \bibverse{John}{5}{28-29} Szenen, in denen er von zwei Seiten im Gericht spricht -
	aber eine der Seiten existiert überhaupt nicht, weil niemand gerecht ist?
	Ist es nicht wahrscheinlicher, dass in poetischem Material der Sprecher eine erhöhte Sicht seiner Unwürdigkeit vor Gott wiedergibt
	und Gott aus Gnade einen Maßstab an Menschen setzt, der tatsächlich realistisch erfüllbar ist?
	Die Tatsache dass das schwarz-weiße Aufteilen der Menschen in Gerechte und Ungerechte im letzten Gericht so häufig formuliert
	wird macht dieses Verständnis deutlich einfacher.
}
In den meisten zitierten Psalmen schreit der Sprecher zu Gott einzugreifen, indem er die ungerechten und gottlosen Unterdrücker bestraft.
Gleichzeitig ist explizit, dass der Sprecher sich selbst auf der Gerechten Seite der Medaillie sieht.
In allen diesen Texten ist klar, dass nicht jeder einzelne Mensch ungerecht im Sinne des Psalmes ist.\footnote{
	Vergleiche \bibverse{Psalms}{14}{1-3} mit v5; \bibverse{Psalms}{5}{10} mit v13; \bibverse{Psalms}{140}{4} mit v14; \bibverse{Psalms}{36}{8} mit v11.
	In \bibverse{Psalms}{10}{} spricht der `Elende' und bittet Gott um Rettung vor dem Gottlosen.
	Die Implikation ist, dass der Elende im Sinne des Psalmes
	nicht ebenso Gottlos ist.
	In \bibverse{Proverbs}{1}{16} nutzt der Lehrer die Aussage, um seinen Schüler von den Ungerechten fern zu halten. Die Aussage ist wiederum nicht,
	dass jeder Mensch in der selben Ungerechten Kategorie ist.
}

Zweitens: Paulus argumentiert nicht, dass `alle Sünden getan haben', sondern dass `alle unter der Sünde sind' - selbst die Juden.
Genau das zeigen die Zitate aus dem Alten Testament wirklich: jeder Mensch, Heide oder Jude, lebt als Mensch unter der Sünde.
Jeder einzelne Mensch sündigt tatsächlich, und niemand kann Gottes Perfektion entsprechen.
Genau das schreibt auch Paulus selbst: ``Denn wir haben sowohl Juden als auch Griechen vorher beschuldigt, alle unter der Sünde zu sein'' (3:9).
Paulus will nicht beweisen, dass `alle Menschen aus Sünde an Gott schuldig geworden sind', oder dass `jeder einzelne Mensch vor Gottes Gericht ungerecht ist'.
Das Problem, unter das Paulus hier Juden ebenso wie die Griechen einschließt ist, unter der Kraft der Sünde versklavt zu sein.\footnote{
	Ein Punkt, den Paulus immer wieder verwendet, beispielsweise in \bibverse{Romans}{6}{17}.
}

Drittens: Es ist in Kapitel 2 bereits klar, dass das Herz einer Person zu allerlei Sünden führt:
``Nach deiner Störrigkeit und deinem unbußfertigen Herzen aber häufst du dir selbst Zorn auf für den Tag des Zorns''.
Das ursprüngliche Problem ist das Herz, die Störrigkeit des Menschen, sein Wesen selbst -
nicht bloß die Tatsache, dass Menschen Schuldig sind und es eine Lösung braucht, damit Gott Sündern die Umkehren wollen vergeben kann.
Es ist die Natur des Menschen an sich in Frage.

Viertens: Es ist auch hilfreich zu verstehen, wie Paulus anderswo mit diesem Problem umgeht:
``denn alle haben gesündigt und es fehlt ihnen die Herrlichkeit Gottes''(3:23) zeigt deutlich,
dass das Problem größer ist als Schuld. Etwas wirkliches, Anfassbares fehlt dem Menschen in seinem sündigen Zustand.

Fünftens: Die Übersetzung `dem Gericht Gottes verfallen' in v19 ist nicht notwendig.
Das griechische ὑπόδικος τῷ θεῷ bezieht sich nicht explizit auf Gottes Gericht, selbst wenn die Elberfelder so übersetzt.\footnote{
	Und klassische Interpretationen wie \Cite{nicnt-romans} lesen hier natürlich eine Szene im Gerichtssaal.
}
Alle Menschen sind `haftbar' vor Gott, allen Menschen - selbst den Juden - ist der Mund gestopft, denn alle Menschen sind in Sünde.
Ob diese Tatsache im Kontext des letzten Gerichts zu verstehen ist oder nicht wird durch diesen Vers nicht markiert.

Zusammenfassend ergeben diese Punkte eine Szene, in der der allgemeine Zustand der Menschheit - der Heiden und Juden - in Frage steht.
Alle Menschen haben ein Herz-Problem\footnote{
	Diese Bezeichnung geht auf \Cite[179]{paul-a-new-covenant-jew} zurück.
	Ich gehe zwar in meiner linguistischen Analyse der δικαι-Wörter anders vor, stimme aber ihren theologischen Argumenten im Großen und Ganzen zu.
}, und dieses Problem muss gelöst werden.
Da keine individuelle Situation im Kontext steht, und Paulus eindeutig über den allgemeinen Zustand aller Menschen unter der Sünde redet
bleibt die Skala für δικαιοω in 3:20 offen und die Gerechtigkeit die das Gesetz nicht erzeugen kann ist die (echte, ontologische) Gerechtigkeit des Herzens.
``Menschen werden nicht aus den Werken des Gesetzes vor Gott gerecht gemacht, \textit{denn durch das Gesetz ist Erkenntnis der Sünde}.'' (3:20).
Das Gesetz erfüllt also nicht die Funktion, Menschen gerecht zu machen, denn stattdessen führt es nur zur Erkenntnis der Sünde.
Das Gesetz ist ein Hilfsmittel, ein Maßstab, an dem Gutes und Böses handeln gemessen werden kann.
Wer diese Gesetze tut wird leben (2:13), aber alle Menschen sind der Sünde verfallen (3:9-18),
und das Gesetz selbst kann den Menschen nicht dazu befähigen so zu handeln (3:20).
Paulus Aussage ist, dass das Gesetz kraftlos ist, Menschen gerecht zu machen.
Es zeigt Sünde auf, kann aber das eigentliche Problem nicht lösen.

Die Verse 2:1 - 3:20 bilden gewissermaßen ein Crux Interpretum und ein Gitter durch das hindurch der Rest des Römerbriefes gelesen wird.
Im Folgenden sehen wir, wie sich unsere Entscheidung in diesen Stellen im Rest des Briefes niederschlägt:
Wir haben ein klares Problem - Menschen sind unter Sünde, und auch das Gesetz ist nicht in der Lage, sie aus dieser Lage heraus gerecht zu machen.
Wenn unsere Interpretation hier korrekt ist sollte dieses Zentrale Problem immer wieder der Kern sein, um den sich Paulus weitere Argumente und seine Lösung dreht.

\subsection{Römer 3:21-26}
\begin{quote}
	Nun aber wurde abseits der Torah Gottes Gerechtigkeit offenbart, bezeugt von der Torah und den Propheten.
	Die Gerechtigkeit Gottes durch den Jesus-Christus-Glauben für alle die glauben.
	Es ist kein Unterschied: Denn alle haben gesündigt, und es fehlt ihnen die Herrlichkeit Gottes.
	Als Geschenk werden sie gerecht gemacht durch seine Gnade, durch den Loskauf in Christus Jesus.
	Ihn stellte Gott in seinem Blut öffentlich hin als Ort der Offenbarung durch Glauben
%	\footnote{
%		Wie genau diese Phrase zu verstehen ist, wird nicht direkt offensichtlich: ist (a) ιλαστηριον durch Glaube zugänglich,
%		oder ist (b) Jesu Treue zu Gottes Plan die Art und Weise, in der diese Offenbarung dargestellt wird?
%		Ich favorisiere die zweite Variante, aber da diese Phrase für unsere theologische Frage nicht relevant ist
%		habe ich hier eine sich nicht festlegende Übersetzung gewählt.
%	}
	zum Beweis seiner Gerechtigkeit im Blick auf das unbestraft lassen der früher geschehenen Sünden
	in der Nachsicht Gottes;
	zum Beweis seiner Gerechtigkeit in der jetzigen Zeit
	sodass er gerecht ist und aus dem Jesus-Glauben gerecht macht.
	\attrib{\bibverse{Romans}{3}{21-26}}
\end{quote}

Einige Kommentare zur Übersetzung an sich sind als aller erstes nötig.
\textbf{Torah} Das ``Gesetz'' ist im ganzen Kontext immer wieder das mosaische Gesetz. Dieses Torah ist es, die die Nationen von Natur aus nicht haben.
		Paulus Argument, gerade in Kapitel zwei, dreht sich immer wieder darum, dass das haben des Gesetzes nicht zum leben führt, sondern das tun.
		Ich sehe im Kontext des Römerbriefes keinen Grund, dieses Gesetz auf die ethnischen Grenzmarker zu beschrenken.

\textbf{Jesus-Christus-Glauben} Die Übersetzung der Phrase πίστεως Ἰησοῦ Χριστοῦ ist äußerst umstritten.
		Die klassischen Positionen sind (a) `der Glauben an Jesus Christus', die sogenannte objektive Sicht, oder (b) `die Treue Jesu Christus', die sogenannte Subjektive Sicht.
		Zu häufig wird die Diskussion über die Bedeutung einer griechischen Phrase als eine theologische Diskussion geführt.\footnote{
			Kontra \Cite[53]{porter-pitts} ``We cannot hope to solve such a delicate and sensitive issue as the meaning of the phrase πίστις Χριστοῦ simply by means of Greek linguistics,
			because there is more at stake – including exegesis and theology''.
			Aber wie kann unsere Exegese und Theologie begründet sein, wenn linguistische Fragen nicht linguistisch geklärt werden?
			Wenn Menschen sprechen und Authoren schreiben meinen sie etwas.
			Unsere Theologie in eine linguistische Frage zurück zu lesen ist eine sehr gefährliche Präzedenz.
		}
		Stattdessen ist es meine Position, dass die detaillierte linguistische Analyse primär gegenüber der Exegese oder gar der systematischen Theologie sein muss.
		Was ist also die beste Interpretation von πιστις Χρισοτου?
		Um diesen Artikel nicht um eine weitere komplexe Diskussion zu erweitern gebe ich hier nur die wichtigsten Argumente aus
		\Cite{grasso-pistis-christou} wieder.
		πιστις ist ein deverbales Substantiv. Solche Substantive können die Argumentstruktur des zugehörigen Verbes behalten.
		Die objektive Sicht bedeutet, dass das direkte Objekt des Verbes zum Genitivkonstrukt des Substantives wird.
		Wenn das Verb πιστευω allerdings ein direktes Objekt im Akkusativ annimmt, dann ist die Bedeutung immer `Glauben an eine Proposition'\footnote{
			\Cite[4.1.1]{grasso-pistis-christou}, cf. \bibverse{John}{11}{26}
		}.
		Da ein objektiver Genitiv diese Argumentstruktur beibehält, muss die Phrase dann `der Glauben, dass (eine unspezifizierte Präposition über) Christus wahr ist' bedeuten.
		Die (notwendige) Bedeutung `Jesus vertrauen' ist in dieser grammatischen Form unmöglich.
		Die subjektive Sicht ist ausgeschlossen, weil (a) das direkte Objekt (oder das indirekte Dativ-Objekt) direkt aus dem Kontext ersichtlich sein muss -
		aber das ist in den wichtigsten Stellen in denen die Phrase auftaucht nicht der Fall (b) es gibt keine Belege dafür, dass griechisch sprechende Kirchenväter
		die Phrase so verstanden haben.
		Die bessere Interpretation (der sogenannte `dritte Weg') ist der in der Übersetzung wiedergegebene Jesus-Christus-Glauben.
		Das ist das Glaubenssystem, dass Jesus Christus als seinen Inhalt hat.
		Eine parallele Formulierung hierzu finden wir in \bibverse{Philippians}{1}{2} συναθλοῦντες τῇ πίστει τοῦ εὐαγγελίου `zusammen für den Glauben des Evangeliums kämpfen'.
		Der Glaube hat das Evangelium als Inhalt - genauso hat der Jesus-Christus-Glauben Jesus Christus als Inhalt.

\textbf{gerecht machen} Es ist Sünde und fehlende Herrlichkeit im Blick, nicht `Übertretung in einer speziellen Situation' oder `Übertretung im Blick des letzten Gerichts'.
		Auch haben wir in den letzten Kapiteln bereits gesehen, wie immer wieder der Charakter der Person, nicht bloß die Sünde-als-einzelne-Tat, oder Sünde-als-Übertretung
		Paulus Anliegen ist.
		Die Skala bleibt also offen, und dementsprechend muss δικαιοω offen gelesen werden.

Der \textbf{Ort der Offenbarung} führt uns zu einer komplizierteren Diskussion. Ich vertrete hier insbesondere die Meinung von Porter in \Cite{mercy-seat-revelation}\footnote{
	Ähnlich wie Bailey in \Cite{jesus-as-the-mercy-seat}.
}.
Das griechische ἱλαστήριον wird in den folgenden Bedeutungen verwendet: (1) Der Deckel der Bundeslade (a) als Ort der Offenbarung\footnote{
	Das ist meine Interpretation. \bibverse{Exodus}{25}{17-22} (das ist die Bedeutung direkt bei der Einsetzung),
	\bibverse{Leviticus}{16}{2}, \bibverse{Numbers}{7}{89}.
	LXX \bibverse{Habakuk}{3}{2} ist eine klare Anspielung auf die Bundeslade (ἐν μέσῳ δύο ζῴων) und beschreibt durchweg eine Offenbarung Gottes an diesem Ort.
}
(b) als Gegenstand, der an Yom Kippur gereinigt wird\bibfoot{Leviticus}{16}{13-15},
(2) als allgemeiner Ort, an dem Opfer gebracht werden\footnote{
	\label{fn:hilstareion-in-iv-mac}
	\bibverse{Amos}{9}{1} muss sich auf die Obere Fläche eines Altares beziehen.
	\bibverse{Ezekiel}{43}{14-20}, wo es mehrere ἱλαστήρια gibt auf denen allerlei Opfer gebracht werden.
	Als letztes auch \bibverse{IVMaccabbees}{17}{22}.
	Diese Stelle wird oft als (die einzige) Stelle zitiert, in der ιλαστηριον die selbe Bedeutung wie in \bibverse{Romans}{3}{25}
	hat, nämlich metonymisch als Opfer, das die Sünden des Volkes wegnimmt.
	Siehe \Cite{mercy-seat-revelation} für eine detailliertere Auseinandersetzung:
	Der ιλαστηριον wird hier verwendet \textit{als Utensil}, mit dem die Brüder zu Tode gefoltert werden.
	Siehe auch \Cite[2,4]{ignatius-of-antioch-to-the-romans}, wo die wilden Tiere der Altar(θυσιαστήριον) sind, auf dem Ignatius geopfert wird (ἵνα διὰ τῶν ὀργάνων τούτων Θεοῦ θυσία εὑρεθῶ).
	Diese Stelle ist in der Formulierung sehr ähnlich zu \bibverse{IVMaccabbees}{17}{22}: διὰ [...] καὶ τοῦ ἱλαστηρίου τοῦ θανάτου αὐτῶν.
}
(3) als Weihgabe in der Heidnischen Welt.\footnote{
	Gemeint sind hier insbesondere Steelen, die als Versöhnungsgeschenk aufgestellt werden.
	Ein Beispiel ist \Cite[16.182]{antiquitates}.
}
Unsere Stelle im Römerbrief ist durchzogen von Wörtern der Offenbarung Gottes:
seine Gerechtigkeit ist offenbart, Jesus ist öffentlich hingestellt, zum Beiweis Gottes Gerechtigkeit.
Wir haben also direkt eine gute Korrelation zwischen Jesus-als-Offenbarung und dem Kontext.
Erzwungen wird diese Interpretation aber letztlich dadurch, dass alle anderen Vorschläge ausgeschlossen sind:
ιλαστηριον bedeutet einfach in keinem extenten Text `Sühneopfer', `Sühneort', `Sühnezeichen' oder ähnliches.
Es gibt \textit{nicht einen Text}, in dem diese Bedeutung vorkommt.\footnote{
	Wie in Fußnote \ref{fn:hilstareion-in-iv-mac} bereits angeschnitten ist die Situation in \bibverse{IVMaccabbees}{17}{22} ähnlich wie in \bibverse{Romans}{3}{25}.
	Es gibt eine bessere Interpretation, die zur Stelle passt und ihre Bedeutung nicht metonymisch (Deckel der Bundeslade, reinterpretiert als Sühneopfer) erzeugen muss.
}
Eine Übersetzung, in der der Deckel der Bundeslade metonymisch für ein Sühneopfer verwendet wird ist äußerst unplausibel,
einfach weil deutlich einfachere Interpretationen existieren.
Der Deckel der Bundeslade selbst ist gerade der Gegenstand, der in \bibverse{Leviticus}{16}{} gereinigt werden muss,
also ist auch eine Interpretation im Sinne von (1.b) nicht sinnvoll.
Die Bundeslade ist kein Altar im eigentlichen Sinne - sie ist der Wohnort Gottes, der von Unreinheit und Sünde befleckt wird.
Das Ziel des Dekontaminationstags ist es gerade, diesen Wohnort Gottes für seine Gegenwart zu säubern, wie wir in Abschnitt \ref{sec:yom-kippur} gesehen haben.
Was auch immer Paulus ausdrücken will, es ist mit Sicherheit nicht, dass Jesus der Ort ist der gereinigt werden muss.\footnote{
	Ein Argument, das selbst Porter nicht angibt, auch wenn wir für die selbe Position argumentieren.
}
Option (2) fällt aus einem ähnlichem Grund aus. Wenn Jesus der Altar ist, was wird dann geopfert? Jesus, also der Altar, selbst?
Option (3) ist in der aktuellen Wissenschaft recht unbeliebt, aber tatsächlich deutlich einfacher zu halten als (1a) oder (2).
Aber Paulus versucht hier ein größeres Problem zu lösen, als dass Menschen und Gott im Streit sind.
Was Jesus tut sollte ein Sünden-Problem lösen.\footnote{
	Für eine moderne Formulierung dieser Sicht (die ich am Ende doch ablehne), siehe \Cite[259ff]{lamb-of-the-free}.
}

Wenn wir ιλαστηριον in 3:25 als Deckel der Bundeslade \textit{als Ort der Offenbarung} (1) lesen, wie erklärt es dann die Passage im ganzen?
Zunächst sollten wir uns fragen, was JHWH im Alten Bund an diesem Ort offenbart hat.
Wenn Mose im Zelt der Begegnung mit Gott spricht, dann geschieht das an der Bundeslade\bibfoot{Numbers}{7}{89}.
Zu diesem Zeitpunkt ist der Großteil des Gesetzes an sich bereits offenbart - aber Mose erhält in seinem Gespräch mit Gott
\textit{Interpretationen und Anwendungen} des Gesetzes auf verschiedene Situationen.
In \bibverse{Numbers}{9}{} besteht das Problem, dass einige Israeliten Leichen-Unrein sind, während das Passah gefeiert werden sollte.
Wie ist mit dieser Situation umzugehen? Können Unreine in dieser speziellen Situation doch am Opfer und Fest teilnehmen?
Ab Vers 9 gibt Jahwe Mose eine positive Antwort: eine Sonderregelung, mit der unter anderem Leichen-Unreine Personen doch am Passah teilhaben können.
In \bibverse{Numbers}{27}{1-11} lesen wir von einem weiteren Sonderfall: Ein Mann stirbt ohne männlichen Erben.
Dürfen seine Töchter stattdessen das Erbe antreten?
Nach diesem Gespräch vor dem Zelt der Begegnung (v2) bringt Mose die Frage vor Gott (v5) und die Antwort ist eine gnädige Sonderfallregelung,
damit die Töchter erben können.

Wenn wir nun zum Römerbrief zurückkehren sehen wir eine sehr ähnliche Funktion, die Jesus Tod erzeugt.
In Christus (durch den Jesus-Christus-Glauben) gibt es eine neue Offenbarung des Gesetzes, die zum Leben führt:
Der wahre Kern, das echte Herz in Gottes Gesetz findet in Jesus Sieg über die Sünde am Kreuz seinen Ausdruck.
Nicht länger ist das Gesetz nur, was Erkenntnis der Sünde bringt (v20), sondern durch die Befähigung in Christus
werden Menschen aus der Sünde erlöst und Tuer des Gesetzes.
Es sind nicht Werke des Gesetzes, sondern der Jesus-Christus-Glauben, der den Menschen befähigt und freisetzt der zu Christus treu ist.
Gleichzeitig rechtfertigt diese Offenbarung auch Gottes eigene Gerechtigkeit:
Wie kann es sein, dass Gott Menschen wie David vergeben hat, wenn doch der Leben hat, der die Torah erfüllt\bibfoot{Leviticus}{18}{5}?
In Christus, dem Ort der neuen Offenbarung, demonstriert Gott, wofür Paulus immer wieder \bibverse{Habakuk}{2}{4} zitiert:
Der durch den Glauben gerechte wird leben!
Diese Botschaft war immer schon wahr,\footnote{
	Das ist kurz gesagt das Argument in Kapitel 4 des Römerbriefes:
	schon bei David und Abraham ist es Glaube, der Sünden vergibt und der Glaube ist es, der dem Menschen als Gerechtigkeit angerechnet wird.
	Wir könnten also sagen: Glaube / Treue = Gerechtigkeit in Gottes Augen; so macht der Glaube Gerecht, weil die beiden Eigenschaften untrennbar zusammen hängen.
	Loyaler Glauben zu Christus ist äquivalent dazu, gerecht zu sein.
	Wer Christus folgt ist automatisch gerecht, und ein Tuer der Torah.
}
aber jetzt durch den Messias gerechtfertigt und allen Menschen offenbart.
Es war nie die Torah, die einen Menschen dazu in die Lage versetzt, gerecht zu sein - es war Glaube (Römer 4).
Und dass diese Tatsache nicht plötzlich Gottes Plan verändert beweist sich darin, wie in Christus jetzt selbst Heiden Tuer der Torah werden (2:14-15).
Das ist das Herz Gottes, die frisch offenbarte Wahrheit in Christus: Gott gibt nicht nur gerechtes Gesetz und setzt einen gerechten Standard (2).
Er selbst greift in seiner selbst aufopfernden Liebe ein, geht über Sünden hinweg, rettet uns aus der Macht der Sünde(3:24b), unter der Juden wie Heiden standen (3:1-20),
und transformiert unser krankes Herz (3:24a).
Der Christus-Glaube, nicht die Werke der Torah, sind es, die einen Menschen gerecht machen.
Und was sagt dieser Christus-Glaube?
Er sagt: alle, die an Gott glauben, gehören zu ihm.
Genau diese Tatsache beweist Gottes Gerechtigkeit in seiner Vergebung für die Sünden in früheren Zeiten.
Schon immer war Gottes Herz Liebe und Vergebung für alle, die ihm folgen - jetzt in Christus endgültig demonstriert und endlich
verbunden mit einer echt wirksamen Transformation des Herzens, die die Werke der Torah nie vollbringen konnten.\footnote{
	Wir denken hier natürlich auch an \bibverse{Ezekiel}{36}{25-27}.
}

\subsection{Römer 5:9-10}
Rom 5:9-10
- Gerecht Gemacht durch Jesus Blut = Transformation = Versöhnung
- Gerettet vor dem Gericht = Gerettet, weil wir tatsächlich durch das Leben in Christus gerecht sind und im Gericht als gerecht befunden werden

\subsection{Römer 5-6 - Der Ungerechtigkeit gestorben, der Gerechtigkeit auferstanden}
tl;dr:
In Römer 5 ab v12 beschreibt Paulus, wie durch Adams Sünde der Tod (=die Verdammnis) in die Welt kam.
Durch Jesus gerechte Tat aber kommt Leben (=Gerecht-Gemacht-Sein; v19).
In Römer 6 beschreibt denselben Kontrast:
Wir sind in Christus gestorben, um damit von der Sünde in der wie gelebt haben gerecht zu sein (v7).
Jetzt könnnen wir in Christus leben. Das passiert, indem wir täglich der Sünde absagen und Werkezeuge der Gerechtigkeit sind, nicht der Ungerechtigkeit.

\subsection{Römer 8:3 - Gottes Urteil über die Sünde}
tl;dr:
Es wird nicht `Jesus für unsere Sünde' verurteilt, sondern `die Sünde im Fleisch'.
Das Ergebnis ist ganz klar, dass die `Rechtsforderung des Gesetzes in uns erfüllt ist', indem wir jetzt durch Loyalität zu Jesus Gerecht handeln -
dadurch, dass wir im Geist leben und die Handlungen des Leibes töten (\bibverse{Romans}{8}{13}).
Es geht nicht darum, wie schuldige Menschen von Gott freigesprochen werden können. Es geht darum, wie die Sünde vernichtet werden kann,
deren Kraft wir Menschen verfallen sind.

